\label{chap:ht-bulk-boundary-line-frontier}

% Theorem environments inherited from main.tex; local macros only.
\providecommand{\kk}{\mathbb{k}}
\providecommand{\cO}{\mathcal O}
\providecommand{\SCop}{\mathsf{SC}^{\mathrm{ch,top}}}
\providecommand{\gSC}{\mathfrak{g}^{\mathrm{SC}}}
\providecommand{\cC}{\mathcal C}
\providecommand{\cL}{\mathcal L}
\providecommand{\cM}{\mathcal M}
\providecommand{\cJ}{\mathcal J}
\providecommand{\Ainf}{\mathsf{A}_{\infty}}
\providecommand{\Linf}{L_\infty}
\providecommand{\Sym}{\operatorname{Sym}}
\providecommand{\PV}{\operatorname{PV}}
\providecommand{\ot}{\otimes}
\providecommand{\CE}{\operatorname{CE}}
\providecommand{\HH}{\operatorname{HH}}
\providecommand{\RHom}{\operatorname{RHom}}
\providecommand{\BBar}{\operatorname{Bar}}
\providecommand{\Cobar}{\operatorname{\Omega}}
\providecommand{\Spec}{\operatorname{Spec}}
\providecommand{\Spf}{\operatorname{Spf}}

\providecommand{\End}{\operatorname{End}}
\providecommand{\Hom}{\operatorname{Hom}}
\providecommand{\Der}{\operatorname{Der}}
\providecommand{\gr}{\operatorname{gr}}
\providecommand{\im}{\operatorname{im}}
\providecommand{\coker}{\operatorname{coker}}
\providecommand{\id}{\operatorname{id}}
\providecommand{\Perf}{\operatorname{Perf}}
\providecommand{\Cl}{\operatorname{Cl}}
\providecommand{\dCrit}{\operatorname{dCrit}}
\providecommand{\Zder}{\operatorname{Z}_{\mathrm{der}}}
\providecommand{\mc}{\operatorname{MC}}
\providecommand{\op}{\mathrm{op}}
\providecommand{\der}{\mathrm{der}}
\providecommand{\Abulk}{A_{\mathrm{bulk}}}
\providecommand{\Bbound}{B_{\partial}}
\providecommand{\Kline}{K_{\mathrm{line}}}
\providecommand{\intr}{\mathrm{intr}}
\providecommand{\free}{\mathrm{free}}
\providecommand{\eff}{\mathrm{eff}}
\providecommand{\st}{\mathrm{st}}
\providecommand{\ex}{\mathrm{ex}}
\providecommand{\minmod}{\mathrm{min}}
\providecommand{\expmod}{\mathrm{exp}}
\providecommand{\hkR}{\mathrm{HKR}}
\providecommand{\sn}{\mathrm{SN}}
\providecommand{\lines}{\mathsf{Lines}}
\providecommand{\hSym}{\widehat{\Sym}}

\section{Toward the global theorem: rich boundaries, absolute lines, and a Steinberg-style endpoint}\label{sec:global}

The exact local theorem package is now in place.  What remains is the global theorem.  This section explains how the local constructions fit into the global program, what is already formal, and what geometric endpoint should ultimately present the full line algebra.

\subsection{The global theorem as a two-comparison problem}

The corrected global triangle \eqref{eq:global-corrected-triangle} will follow once two comparison maps are supplied:

\begin{enumerate}[label=(\arabic*),leftmargin=1.8em]
\item a \emph{bulk/boundary comparison}
\[
\beta_{\der}:\Abulk\longrightarrow \Zder(B_b)
\]
for a rich boundary condition $b$;
\item a \emph{boundary/line comparison}
\[
\End_{\cC_{\mathrm{line}}}(M_b)\simeq B_b
\]
showing that the chosen boundary vacuum $M_b$ generates the line category.
\end{enumerate}

Once these exist, Proposition~\ref{prop:formal-global-triangle} finishes the job.  So the theorem is not one enormous mystery.  It is two comparison theorems plus Morita invariance.

\subsection{Why the local exact sector matters for the global problem}

The local theorem package proved above gives a template for the global attack:
\begin{itemize}[leftmargin=1.8em]
\item \textbf{intrinsic boundary algebra}: the open sector should be defined without arbitrary splittings;
\item \textbf{exact local line models}: local line algebras should arise from a canonical coalgebraic or Koszul-dual presentation;
\item \textbf{Kuranishi reduction}: linearly massive modes should be integrated out exactly;
\item \textbf{derived center}: bulk should be recovered from lines by Hochschild cochains;
\item \textbf{stabilization calculus}: universal contractible and Morita-trivial blocks must be removed before comparing theories.
\end{itemize}
The local exact sector shows that each of these moves is real, not aspirational.

\subsection{The absolute line algebra versus boundary-adapted local line algebras}

It is essential not to confuse two different line objects:
\begin{enumerate}[label=(\roman*),leftmargin=1.8em]
\item the \emph{absolute} line algebra $A^!$ of the global theory, conjecturally presenting all line operators and appearing in \cite{DNP25};
\item the \emph{boundary-adapted pointed} line algebras $K_{F,p}$ and $K_\kappa$ constructed in Sections~\ref{sec:exact-lines}--\ref{sec:classification}.
\end{enumerate}
The latter are local exact models attached to a chosen boundary vacuum.  They should be viewed as boundary charts on the absolute line category.  They are not globally equal to $A^!$, but they are the correct local testing ground in which the shape of the global theorem becomes rigid.

\subsection{The Steinberg presentation theorem}

The phrase ``Chriss--Ginzburgification'' is not decorative.
It names a methodological demand: a mature representation-theoretic
algebra should eventually be presented by correspondences.
The Steinberg variety $Z = \widetilde{\cN} \times_{\mathfrak g}
\widetilde{\cN}$ is a correspondence between two copies of the
Springer resolution; its convolution algebra $H_*(Z)$ recovers the
Hecke algebra $\cH_q(W)$.  The bar complex $\barB(\cA)$ on
$\FM_k(\C) \times \Conf_k^{<}(\R)$ is a correspondence between
holomorphic collision data and topological ordering data; its
convolution algebra recovers the Swiss-cheese composition.
The following theorem makes this analogy precise.

\begin{theorem}[Algebraic Steinberg presentation; \ClaimStatusProvedHere]
\label{thm:steinberg-presentation}
\index{Steinberg presentation|textbf}
\index{Swiss-cheese!Steinberg presentation|textbf}
Let $\cA$ be a chiral algebra on a smooth curve~$X$ satisfying the
standing hypotheses, and let $T = (\Bbound, \cC_{\mathrm{line}})$
be the associated $\SCop$-algebra
\textup{(}Theorem~\textup{\ref{thm:bar-swiss-cheese}}\textup{)}.
Then the Swiss-cheese convolution $L_\infty$-algebra $\gSC_T$
\textup{(}Definition~\textup{\ref{def:sc-convolution}}\textup{)}
\textbf{presents} the bulk/boundary/line Koszul triangle in the
following sense.
\begin{enumerate}[label=\textup{(\roman*)}]
\item \textbf{Presentation.}
  The operadic composition maps of $\SCop$ are recovered from the
  bar differential and coproduct on $\barB(\cA)$:
  \begin{align*}
  \text{closed composition}
    &= d_{\barB}
    &&\text{\textup{(}OPE residues on $\FM_k(X)$\textup{)}},\\
  \text{open composition}
    &= \Delta
    &&\text{\textup{(}deconcatenation on $\Conf_k^{<}(\R)$\textup{)}},\\
  \text{mixed composition}
    &= d_{\barB}\circ\Delta
      = (d_{\barB}\otimes\id + \id\otimes d_{\barB})\circ\Delta
    &&\text{\textup{(}coderivation\textup{)}}.
  \end{align*}
  The no-open-to-closed rule is the directionality of the
  correspondence: $\Delta$ acts on bar degree and does not produce
  new OPE poles.  These three compositions, together with the
  canonical MC element $\alpha_T$ of
  Theorem~\textup{\ref{thm:3d-universal-mc}}, determine the full
  $\SCop$-algebra structure by homotopy transfer
  \textup{(}Theorem~\textup{\ref{thm:homotopy-Koszul}}\textup{)}.

\item \textbf{Line algebra from convolution.}
  The module category of the Swiss-cheese convolution algebra is
  equivalent to the Koszul dual module category:
  \[
  \cC_{\mathrm{line}}
  \;\simeq\;
  \mathrm{Rep}(\gSC_T)
  \;\simeq\;
  A^!\textup{-mod}.
  \]
  The first equivalence is the definition of the open-colour sector
  of $\alpha_T$
  \textup{(}Proposition~\textup{\ref{prop:alpha-projections}(ii)}\textup{)};
  the second is the unconditional line-operator theorem
  \textup{(}Theorem~\textup{\ref{thm:lines_as_modules}}\textup{)}.

\item \textbf{Bulk from derived center.}
  The derived center of the convolution algebra recovers the bulk:
  \[
  \Abulk
  \;\simeq\;
  \Zder(\gSC_T)
  \;\simeq\;
  \HH^\bullet_{\mathrm{ch}}(A^!).
  \]
  The first equivalence follows from the $(1,1)$-projection of
  $\alpha_T$
  \textup{(}Proposition~\textup{\ref{prop:alpha-projections}(iv)}\textup{)};
  the second is the Hochschild bridge
  \textup{(}Theorem~\textup{\ref{thm:bulk_hochschild}}\textup{)}.

\item \textbf{Braiding from monodromy.}
  The spectral $R$-matrix $R(z)$ is the monodromy of the
  correspondence: the $(2,0)$-projection of $\alpha_T$ restricted to
  the ordered collision stratum of $\FM_2(\C)$
  \textup{(}Proposition~\textup{\ref{prop:alpha-projections}(iii)}\textup{)}.
  The Yang--Baxter equation is the MC equation at arity $(3,0)$.

\item \textbf{Genus tower from curvature.}
  At genus $g \ge 1$, the self-loop restriction of $\alpha_T$
  recovers the modular characteristic $\kappa(\Bbound)\cdot\omega_g$
  and the shadow Postnikov tower
  \textup{(}Proposition~\textup{\ref{prop:alpha-projections}(v)}\textup{)}.
\end{enumerate}
\end{theorem}

\begin{proof}
The theorem assembles five proved results into a single package.
Statement~(i) is Theorem~\ref{thm:bar-swiss-cheese} (the bar
complex is an $\SCop$-algebra) together with
Theorem~\ref{thm:homotopy-Koszul} (homotopy-Koszulity ensures
the bar-cobar resolution recovers the operad).
Statement~(ii) combines
Proposition~\ref{prop:alpha-projections}(ii) (the open projection
gives $A^!$-modules) with Theorem~\ref{thm:lines_as_modules}
(now unconditional).
Statement~(iii) combines
Proposition~\ref{prop:alpha-projections}(iv) (the mixed projection
gives the bulk-to-boundary map) with
Theorem~\ref{thm:bulk_hochschild} (bulk as chiral Hochschild).
Statement~(iv) is Proposition~\ref{prop:alpha-projections}(iii)
together with Theorem~\ref{thm:YBE}.
Statement~(v) is Proposition~\ref{prop:alpha-projections}(v)
together with the genus-tower identification from Vol~I
(Theorem~D, bridge MC5).

The key point is that all five statements are projections of a
\emph{single} MC element $\alpha_T \in \mc(\gSC_T)$: the bulk,
boundary, lines, braiding, and genus tower are not independent
structures assembled by hand, but different faces of one
correspondence.  This is the content of the presentation.
\end{proof}

\begin{remark}[The Steinberg dictionary]
\label{rem:steinberg-dictionary}
\index{Steinberg!dictionary}
The correspondence between geometric representation theory
and holomorphic-topological algebra is:
\begin{center}
\renewcommand{\arraystretch}{1.3}
\begin{tabular}{lll}
\hline
\textbf{Steinberg geometry} & \textbf{HT bar complex} & \textbf{Role} \\
\hline
Flag variety $G/B$ & $\FM_k(\C)$ & Collision/resolution data \\
Nilpotent cone $\cN$ & Collision divisors $D_S$ & Singular locus \\
Springer resolution $\widetilde\cN \to \cN$
  & $\FM_k(\C) \to \Conf_k(\C)$ & Compactification \\
Steinberg $Z = \widetilde\cN \times_{\mathfrak g} \widetilde\cN$
  & $\FM_k(\C) \times \Conf_k^{<}(\R)$ & Correspondence \\
Convolution $H_*(Z) = \cH_q(W)$
  & $\gSC_T$ = SC convolution $L_\infty$ & Algebra from corr. \\
Hecke algebra $\cH_q(W)$
  & $A^!$ (Koszul dual/boundary) & Presented algebra \\
$\mathrm{Rep}(\cH_q) = \cO_q$
  & $\cC_{\mathrm{line}} = A^!\text{-mod}$ & Module category \\
Center $Z(\cH_q)$
  & $\Abulk = \Zder(\gSC_T)$ & Derived center \\
$q$-parameter
  & Spectral parameter $z$ & Deformation \\
\hline
\end{tabular}
\end{center}
The two colours of the Swiss-cheese operad are the two factors of
the correspondence.  The bar-cobar adjunction is the
convolution-restriction adjunction.  The no-open-to-closed rule
is the directionality of the Springer correspondence (representations
restrict from $G$ to $B$ but not conversely).
\end{remark}

\subsection{Geometric enhancement}

The algebraic Steinberg presentation
(Theorem~\ref{thm:steinberg-presentation}) is unconditional.
The geometric enhancement---constructing the correspondence space
as a derived self-intersection---remains conjectural.

\begin{conjecture}[Geometric HT Steinberg object; \ClaimStatusConjectured]
\label{conj:geometric-steinberg}
\index{Steinberg!geometric conjecture}
To a sufficiently well-behaved perturbative 3d
holomorphic--topological QFT $T$ there is attached:
\begin{enumerate}[label=\textup{(\alph*)}]
\item a derived formal moduli stack $\cM_{\mathrm{vac}}(T)$ of
  local bulk vacua, equipped with a $(-2)$-shifted symplectic
  structure;
\item for every rich boundary condition $b$, a shifted-Lagrangian
  substack $\cL_b \hookrightarrow \cM_{\mathrm{vac}}(T)$;
\item a derived self-intersection (the HT Steinberg object)
  \[
  \mathfrak S_b
  \;:=\;
  \cL_b \times_{\cM_{\mathrm{vac}}(T)} \cL_b,
  \]
  carrying a $(-1)$-shifted symplectic structure from
  complementary Lagrangians.
\end{enumerate}
The convolution algebra of $\mathfrak S_b$ in Borel--Moore homology
is quasi-isomorphic to the Swiss-cheese convolution $L_\infty$-algebra:
$H_*^{\mathrm{BM}}(\mathfrak S_b) \simeq \gSC_T$.
The algebraic Steinberg presentation
\textup{(}Theorem~\textup{\ref{thm:steinberg-presentation}}\textup{)}
is the shadow of this geometric identification.
\end{conjecture}

Evidence for the geometric enhancement:
\begin{itemize}[leftmargin=1.8em]
\item In the exact LG sector, the boundary algebra is literally
  the algebra of functions on a derived zero locus, and the bulk
  is functions on a shifted cotangent of that zero locus.
\item The pointed line algebra is the cobar of a canonical coalgebra
  presenting the local equations---the algebraic shadow of the
  Jacobi neighbourhood of $\cL_b$ in $\cM_{\mathrm{vac}}(T)$.
\item The stabilization theorems identify the universal contractible
  and Clifford blocks that should correspond to trivial
  correspondences or Morita-trivial kernels.
\item The dg-shifted-Yangian package of \cite{DNP25} has exactly
  the structure of a convolution algebra with coproduct and
  $r$-matrix.
\end{itemize}

%% ----------------------------------------------------------------
\subsection{Toward the geometric Steinberg: the $(-2)$-shifted symplectic foundation}
\label{ssec:geometric-steinberg-foundation}

The Chriss--Ginzburg strategy for the classical Steinberg variety
proceeds in two steps: first establish the geometry of the
correspondence (that the ambient space carries the correct symplectic
structure, that the Lagrangians are isotropic, that the fiber product
is well-defined in the derived sense), then let the convolution
algebra project out by formal properties of Borel--Moore homology.
The same strategy applies here.  We prove ingredient~(a) of
Conjecture~\ref{conj:geometric-steinberg}: the derived moduli of
vacua $\cM_{\mathrm{vac}}(T)$ carries a $(-2)$-shifted symplectic
structure.  This is the foundation on which the Lagrangian boundary
conditions and the Steinberg fiber product rest.

\begin{proposition}[{$(-2)$-shifted symplectic structure: worked examples}]
\label{prop:minus-two-shifted-examples}
\ClaimStatusProvedHere
For each of the three standard HT theories, the derived moduli of
vacua on the formal $3$-disk $D^3 = D^2_z \times D^1_t$ carries a
$(-2)$-shifted symplectic structure $\omega_{-2}$.
\begin{enumerate}[label=\textup{(\roman*)},leftmargin=2em]
\item \textbf{Free $\beta\gamma$.}
  $\cM_{\mathrm{vac}} = T^*[-2]\bigl(\operatorname{Maps}(D^2, V)\bigr)$
  with the canonical Liouville form.

\item \textbf{Landau--Ginzburg.}
  $\cM_{\mathrm{vac}} = \operatorname{Crit}[-2](W)$, the derived
  critical locus of the superpotential $W$, with the $(-2)$-shifted
  structure induced by the Hessian of~$W$.

\item \textbf{Chern--Simons.}
  $\cM_{\mathrm{vac}} = \operatorname{Flat}_G^{\mathrm{der}}(D^3)$,
  the derived stack of flat $G$-bundles on the formal $3$-disk,
  $(-2)$-shifted symplectic by PTVV \cite{PTVV13} applied to the
  Chern--Simons functional $\operatorname{CS} \colon
  \operatorname{Flat}_G^{\mathrm{der}}(D^3) \to \mathbb{A}^1[-2]$.
\end{enumerate}
\end{proposition}

\begin{proof}
(i)\enspace The BV action for the free $\beta\gamma$ system is
quadratic: $S_{\mathrm{BV}} = \int_{D^3} \beta\, \bar\partial\gamma$.
The derived space of solutions to $\delta S = 0$ on the formal disk
is the derived mapping stack $\operatorname{Maps}(D^2, V)$, and the
full BV field--antifield complex is its shifted cotangent bundle
$T^*[-2](\operatorname{Maps}(D^2, V))$.  The BV antibracket is a
degree $(-1)$ Poisson structure on fields, which by
Pantev--To\"en--Vaqui\'e--Vezzosi \cite{PTVV13} is equivalent to a
$(-2)$-shifted symplectic structure on the derived critical locus of
the action.  Since $S_{\mathrm{BV}}$ is quadratic, the critical locus
is smooth and the shifted cotangent carries the canonical Liouville
$(-2)$-form.

(ii)\enspace For Landau--Ginzburg with superpotential $W \colon X \to
\mathbb{A}^1$, the derived critical locus
$\operatorname{Crit}^{\mathrm{der}}(W) = X
\times^{\mathbb{R}}_{X \times X} \Gamma_{dW}$ carries a
$(-1)$-shifted symplectic structure by PTVV \cite{PTVV13},
Theorem~2.1.  The BV theory on $D^3 = D^2_z \times D^1_t$
has field space $\mathcal{F} = \Omega^{0,\bullet}(D^2, X) \otimes
\Omega^\bullet(D^1)$.  The BV antibracket pairs fields and
antifields with a degree $(-1)$ pairing; on the derived critical
locus $\cM_{\mathrm{vac}} =
\operatorname{Crit}^{\mathrm{der}}(S_{\mathrm{BV}})$, the
embedding as a zero-locus of $\delta S_{\mathrm{BV}}$ contributes
an additional $(-1)$ shift via the conormal exact triangle
$N^*[-1] \to \mathbb{L}_{\cM_{\mathrm{vac}}} \to
\mathbb{L}_{\mathcal{F}}|_{\cM_{\mathrm{vac}}}$, yielding total
shift $(-1)+(-1) = (-2)$.  At each critical point of $W$, the
$(-2)$-form is controlled by $\operatorname{Hess}(W)$.

(iii)\enspace The PTVV theorem \cite{PTVV13} applies to
\emph{compact oriented} $n$-manifolds: $\operatorname{Flat}_G^{\mathrm{der}}(M)$
carries a $(2-n)$-shifted symplectic structure.  For the formal
$3$-disk $D^3$, which is contractible and has boundary, the PTVV
theorem does not apply directly.  Instead, the $(-2)$-shifted
structure arises from the BV formalism: the Chern--Simons BV
action $S_{\mathrm{BV}} = \int_{D^3} \operatorname{tr}\bigl(A
\wedge dA + \tfrac{2}{3}A^3 + A^+ \wedge d_A c\bigr)$ defines a
degree~$0$ function on the space of BV fields $\mathcal{F} =
\Omega^\bullet(D^3, \mathfrak{g})[1]$, which carries a degree
$(-1)$ BV antibracket from the invariant pairing on $\mathfrak{g}$
and integration over $D^3$.  The derived critical locus
$\cM_{\mathrm{vac}} = \operatorname{Crit}^{\mathrm{der}}(S_{\mathrm{BV}})$
is the derived stack of flat connections on $D^3$
(which is a point: $\operatorname{Flat}_G^{\mathrm{der}}(D^3)
\simeq B\mathfrak{g}$ formally).  As in case~(ii), the embedding
as a zero-locus shifts by $(-1)$, giving total shift $(-2)$.
Concretely, $\omega_{-2} = dd_{\mathrm{dR}} \operatorname{CS}$
where $\operatorname{CS}(A) = \int_{D^3}
\operatorname{tr}\bigl(A \wedge dA + \tfrac{2}{3} A^3\bigr)$.
\end{proof}

\begin{proposition}[{$(-2)$-shifted symplectic structure: general HT sigma model}]
\label{prop:minus-two-shifted-general}
\ClaimStatusConjectured
Let $T$ be a $3d$ theory on $\mathbb{C}_z \times \mathbb{R}_t$ whose
BV data satisfy~\textup{(H1)}.  Then the derived moduli of vacua
$\cM_{\mathrm{vac}}(T)$ on the formal $3$-disk carries a
$(-2)$-shifted symplectic structure.
\end{proposition}

\begin{proof}[Proof sketch]
The BV formalism assigns to a theory on a manifold $M$ with boundary
$\partial M$ a $(-1)$-shifted symplectic structure on the derived
space of fields---this is the content of the
Cattaneo--Mnev--Reshetikhin (CMR) framework \cite{CMR14}.
Specifically, the BV antibracket on bulk fields gives a degree
$(-1)$ Poisson structure; on the derived critical locus of the BV
action, this is equivalent to $(-1)$-shifted symplectic data.

Axiom~(H1) guarantees a non-degenerate BV pairing
$\langle\!\langle -,- \rangle\!\rangle$ of cohomological degree
$(-1)$ on the space of fields $\mathcal{F}$.  This pairing is the
BV antibracket, making $(\mathcal{F}, S_{\mathrm{BV}})$ a
$(-1)$-shifted symplectic derived stack in the sense of
PTVV \cite{PTVV13}.  The derived critical locus
$\cM_{\mathrm{vac}}(T) =
\operatorname{Crit}^{\mathrm{der}}(S_{\mathrm{BV}}) \subset
\mathcal{F}$ is the derived zero-locus of the section
$\delta S_{\mathrm{BV}} \in \Gamma(\mathbb{L}_{\mathcal{F}})$.
The cotangent complex of this zero-locus fits into the
conormal exact triangle
\[
  N^*_{\cM_{\mathrm{vac}}/\mathcal{F}}[-1]
  \;\longrightarrow\;
  \mathbb{L}_{\cM_{\mathrm{vac}}}
  \;\longrightarrow\;
  \mathbb{L}_{\mathcal{F}}|_{\cM_{\mathrm{vac}}},
\]
where $N^*_{\cM_{\mathrm{vac}}/\mathcal{F}} \simeq
\mathbb{T}_{\mathcal{F}}|_{\cM_{\mathrm{vac}}}$ via the Hessian
of $S_{\mathrm{BV}}$.  The $(-1)$-shifted symplectic form on
$\mathcal{F}$ identifies
$\mathbb{T}_{\mathcal{F}} \simeq \mathbb{L}_{\mathcal{F}}[-1]$;
composing this with the conormal identification gives
$\mathbb{T}_{\cM_{\mathrm{vac}}} \simeq
\mathbb{L}_{\cM_{\mathrm{vac}}}[-2]$, which is a $(-2)$-shifted
symplectic structure.  The additional shift of $(-1)$ beyond the
ambient $(-1)$ is the standard derived-intersection phenomenon:
cutting out by one section shifts the symplectic degree by $(-1)$.

The mechanism is uniform across all three worked examples
(Proposition~\ref{prop:minus-two-shifted-examples}), which verify
it in each geometric incarnation.  Making this argument fully
rigorous for a general theory $T$ requires establishing that the
BV data of~(H1) define a $(-1)$-shifted symplectic structure on
$\mathcal{F}$ in the sense of derived algebraic geometry, which
we expect but do not prove here.
\end{proof}

\begin{remark}[Costello--Gwilliam factorization perspective]
\label{rem:CG-factorization-shifted}
In the Costello--Gwilliam framework \cite{CG17a,CG17b}, the
factorization algebra $\operatorname{Obs}^q(D^3)$ of quantum
observables on the formal disk carries a $P_0$-structure---a
degree~$(-1)$ Poisson bracket coming from the BV antibracket.  The
$(-2)$-shifted symplectic structure on $\cM_{\mathrm{vac}}(T)$
established above is the \emph{classical shadow} of this
$P_0$-structure: the factorization algebra of \emph{classical}
observables $\operatorname{Obs}^{\mathrm{cl}}(D^3)$ is equivalent to
functions on the $(-2)$-shifted symplectic derived stack, and the
$P_0$-bracket is the Poisson bracket associated to $\omega_{-2}$.
This ensures that the geometric data
(Propositions~\ref{prop:minus-two-shifted-examples}
and~\ref{prop:minus-two-shifted-general}) and the algebraic data
(the Swiss-cheese convolution algebra) are two readings of the same
structure---precisely the Chriss--Ginzburg philosophy applied to the
holomorphic-topological setting.
\end{remark}

%% ----------------------------------------------------------------
\subsection{Boundary Lagrangians: the shifted isotropic condition}
\label{ssec:boundary-lagrangians}

We now prove ingredient~(b) of
Conjecture~\ref{conj:geometric-steinberg}: every rich boundary
condition $b$ defines a shifted-Lagrangian
$\cL_b \hookrightarrow \cM_{\mathrm{vac}}(T)$.  We establish this
first in the two nontrivial worked examples, then in general.

\begin{proposition}[Lagrangian boundary conditions: Landau--Ginzburg]
\label{prop:lagrangian-lg}
\ClaimStatusProvedHere
Let $T$ be a Landau--Ginzburg theory with superpotential
$W \colon V \to \mathbb{C}$, and let $b$ be a boundary condition.
Then the derived boundary locus
\[
  \cL_b \;=\; \operatorname{dCrit}(W|_b)
  \;=\; \bigl\{\, x \in V : dW(x) = 0,\;\text{boundary condition } b \,\bigr\}
\]
is Lagrangian in the $(-2)$-shifted symplectic stack
$\cM_{\mathrm{vac}} = \operatorname{Crit}[-2](W)$.
\end{proposition}

\begin{proof}
The boundary condition $b$ cuts out a derived subscheme of the
target $V$; the restriction $W|_b$ defines a function on that
subscheme.  The derived critical locus $\cL_b =
\operatorname{dCrit}(W|_b)$ is then the zero locus of the section
$dW|_b$ of the cotangent bundle $\Omega^1_{V|_b}$.

In the $(-2)$-shifted symplectic geometry of
$\cM_{\mathrm{vac}} = \operatorname{Crit}[-2](W)$, this zero-locus
inherits a Lagrangian structure by the exact-sequence criterion.
Concretely, the conormal sequence for the inclusion
$\cL_b \hookrightarrow \cM_{\mathrm{vac}}$ reads
\[
  \mathbb{L}_{\cL_b / \cM_{\mathrm{vac}}}
  \;\longrightarrow\;
  \mathbb{L}_{\cL_b}
  \;\longrightarrow\;
  \mathbb{L}_{\cM_{\mathrm{vac}}}|_{\cL_b}.
\]
The $(-2)$-shifted symplectic form on $\cM_{\mathrm{vac}}$ provides
an identification
$\mathbb{T}_{\cM_{\mathrm{vac}}} \simeq
\mathbb{L}_{\cM_{\mathrm{vac}}}[-2]$.  Restricting to $\cL_b$ and
comparing with the conormal sequence gives the isotropic condition
$N^*_{\cL_b / \cM_{\mathrm{vac}}} \simeq
\mathbb{T}_{\cL_b}[-1]$.  Since $\cL_b$ is cut out by a section
of a vector bundle (namely $dW|_b$), this isotropic condition is
automatically an equivalence---the zero-locus of a section is
Lagrangian by the PTVV criterion \cite{PTVV13}.  This is the
$(-2)$-shifted analogue of the classical fact that the graph of a
closed $1$-form is Lagrangian in the cotangent bundle.
\end{proof}

\begin{proposition}[Lagrangian boundary conditions: Chern--Simons]
\label{prop:lagrangian-cs}
\ClaimStatusProvedHere
For Chern--Simons theory with gauge group~$G$, let $N \subset G$ be
a maximal unipotent subgroup and consider the Nahm boundary
condition.  Then
\[
  \cL_{\mathrm{Nahm}}
  \;=\;
  \operatorname{Bun}_G^{\mathrm{flat}}
    \bigl(D^2,\; \partial D^2 \to G/N\bigr)
\]
is Lagrangian in the $(-2)$-shifted symplectic stack
$\cM_{\mathrm{vac}} =
\operatorname{Flat}_G^{\mathrm{der}}(D^3)$.
\end{proposition}

\begin{proof}
By PTVV \cite{PTVV13}, restriction of flat bundles from a
$3$-manifold to its boundary is a Lagrangian morphism:
the derived restriction map
\[
  \operatorname{Flat}_G^{\mathrm{der}}(D^3)
  \;\longrightarrow\;
  \operatorname{Flat}_G^{\mathrm{der}}(\partial D^3)
\]
sends the $(-2)$-shifted symplectic structure on the source to a
Lagrangian structure on the target.  More precisely, this is the
content of the Calaque--Haugseng--Scheimbauer theorem
\cite{CHS21}: the restriction functor from a cobordism to its
boundary is a Lagrangian correspondence in the $(\infty,n)$-category
of shifted-symplectic stacks.

The Nahm boundary condition imposes a reduction of structure:
on the boundary $\partial D^2$, the $G$-bundle is
equipped with a reduction to $G/N$.  This defines a derived
substack $\cL_{\mathrm{Nahm}}$ of the boundary moduli.
The Lagrangian property transfers to $\cL_{\mathrm{Nahm}}$ by
the following argument.  The reduction to $G/N$ is the
zero-locus of the section $\mu_N \colon
\operatorname{Flat}_G^{\mathrm{der}}(\partial D^3) \to
\mathfrak{n}^*[-1]$ given by the moment map for the $N$-action.
By the PTVV zero-locus criterion \cite{PTVV13}, a derived
zero-locus of an isotropic section inherits a Lagrangian
structure.  The section $\mu_N$ is isotropic because $N$ is
coisotropic in $G$ with respect to the Killing form (the Lie
algebra $\mathfrak{n}$ is a maximal isotropic subspace of
$\mathfrak{g}$ under the invariant pairing).  At the classical
level, this recovers the Springer resolution
$T^*(G/B) \to \mathcal{N}$ as the Lagrangian correspondence.
\end{proof}

\begin{proposition}[Lagrangian boundary conditions: general HT theory]
\label{prop:lagrangian-general}
\ClaimStatusProvedHere
Let $T$ be a $3d$ theory on $\mathbb{C}_z \times \mathbb{R}_t$
whose BV data satisfy~\textup{(H1)}, and let $b$ be a rich
boundary condition in the sense of
Definition~\textup{\ref{def:rich-boundary}}.  Then $b$ defines a
Lagrangian
$\cL_b \hookrightarrow \cM_{\mathrm{vac}}(T)$ in the
$(-2)$-shifted symplectic stack of
Proposition~\textup{\ref{prop:minus-two-shifted-general}}.
\end{proposition}

\begin{proof}
In the Cattaneo--Mnev--Reshetikhin (CMR) boundary BV framework
\cite{CMR14}, a boundary condition~$b$ on a $3$-manifold $M$ with
boundary $\partial M$ is a choice of Lagrangian in the
$(-1)$-shifted symplectic space of boundary fields
$\mathcal{F}_{\partial M}$.  Passing to the derived critical locus
of the bulk BV action---which carries the
$(-2)$-shifted symplectic structure---the boundary condition
restricts to an isotropic substack
$\cL_b \hookrightarrow \cM_{\mathrm{vac}}(T)$.

The isotropic condition is automatic: the BV boundary condition
requires $\omega_{-2}|_{\cL_b} = 0$.  What requires proof is the
\emph{non-degeneracy} upgrading isotropic to Lagrangian, i.e.\
that the conormal sequence
\[
  N^*_{\cL_b / \cM_{\mathrm{vac}}}
  \;\xrightarrow{\;\sim\;}
  \mathbb{T}_{\cL_b}[-1]
\]
is an equivalence rather than merely a map.  This is precisely where
the ``richness'' condition on~$b$ is used.  A rich boundary condition,
by definition, has the property that the boundary fields carry a
complete set of propagating degrees of freedom---the boundary BV
complex is quasi-isomorphic to the full restriction of the bulk
complex.  In the conormal sequence, this means the cofiber of
$N^*_{\cL_b / \cM_{\mathrm{vac}}} \to \mathbb{T}_{\cL_b}[-1]$
vanishes: the map is an equivalence.

Repackaging in PTVV language: a rich BV boundary condition is a
derived zero-locus of a section of a vector bundle over
$\cM_{\mathrm{vac}}(T)$ (the section being the boundary equations
of motion), and the zero-locus of a section is automatically
Lagrangian by PTVV \cite{PTVV13}, Theorem~2.9.  The two worked
examples (Propositions~\ref{prop:lagrangian-lg}
and~\ref{prop:lagrangian-cs}) verify this general mechanism in
their respective geometric incarnations.
\end{proof}

\begin{remark}[Physical interpretation: half-BPS condition]
\label{rem:lagrangian-half-susy}
The $(-2)$-shifted Lagrangian condition on $\cL_b$ is the
mathematical encoding of ``the boundary condition preserves half
the supersymmetry'' in physics language.  A $(-2)$-shifted
symplectic structure has $2n$ ``fermionic directions'' (in the
sense of the BV--BRST complex); a Lagrangian kills exactly half
of them.  This is the derived-geometric avatar of the classical
half-BPS condition: a brane wrapping a Lagrangian submanifold of
a symplectic target preserves half the supercharges.
\end{remark}

%% ----------------------------------------------------------------
\subsection{The convolution identification: from fiber products to the
Swiss-cheese algebra}
\label{ssec:convolution-identification}

\providecommand{\Steinb}{\mathfrak{S}}
\providecommand{\BM}{\mathrm{BM}}
\providecommand{\Conf}{\operatorname{Conf}}
\providecommand{\FM}{\operatorname{FM}}

We now prove ingredient~(c) of
Conjecture~\ref{conj:geometric-steinberg}: the convolution algebra on
the Steinberg fiber product $\Steinb_b = \cL_b
\times_{\cM_{\mathrm{vac}}} \cL_b$ reproduces the Swiss-cheese
algebra structure.  The key is that iterated fiber products over
$\cM_{\mathrm{vac}}$ decompose as products of holomorphic and
topological configuration spaces, and that convolution on
Borel--Moore homology recovers the bar differential and coproduct
simultaneously.

\begin{construction}[Iterated fiber product]
\label{cons:iterated-fiber}
\index{Steinberg object!iterated fiber product}
For a boundary condition $b$ in a theory $T$ on $\C \times \R_{\ge
0}$, and $k \ge 1$, define the \emph{$k$-fold iterated Steinberg
object}
\begin{equation}\label{eq:iterated-steinberg}
\Steinb_b^{(k)}
\;:=\;
\underbrace{
  \cL_b \times_{\cM_{\mathrm{vac}}}
  \cL_b \times_{\cM_{\mathrm{vac}}}
  \cdots \times_{\cM_{\mathrm{vac}}}
  \cL_b
}_{k+1 \text{ copies of } \cL_b},
\end{equation}
where each fiber product is taken in the derived sense (using the
$(-2)$-shifted symplectic structure of
Proposition~\ref{prop:minus-two-shifted-general}).  Geometrically,
$\Steinb_b^{(k)}$ parametrizes configurations of $k$ ``interaction
points'' at which boundary conditions meet through the bulk vacuum
moduli.

The $k$-fold object admits a factored description over the base of
bulk field configurations:
\begin{equation}\label{eq:steinberg-factored}
\Steinb_b^{(k)}
\;\simeq\;
\FM_k(\C) \times \Conf_k^{<}(\R),
\end{equation}
where $\FM_k(\C)$ is the Fulton--MacPherson compactification of
$k$ points in $\C$ (encoding the holomorphic positions of the
interaction points) and $\Conf_k^{<}(\R) = \{t_1 < \cdots < t_k\}$
is the ordered configuration space of $k$ points on $\R$ (encoding
the topological ordering along the interval).  The equivalence
\eqref{eq:steinberg-factored} is \emph{expected} over the formal
completion at the vacuum, motivated by the product structure of the
ambient space $\C_z \times \R_t$: the fiber product constraint from
$\cM_{\mathrm{vac}}$ acts on the holomorphic direction (matching
bulk fields at each interaction point), while the ordering
constraint acts on the topological direction (the $E_1$-structure
along~$\R$).  A complete proof would require showing that the
derived fiber product over $\cM_{\mathrm{vac}}$ decomposes as a
product of the holomorphic and topological moduli---this holds when
the theory genuinely factorizes as a product of a holomorphic theory
on $\C$ and a topological theory on $\R$, which is the content of
the HT factorization axiom~(H4).
\end{construction}

\begin{proposition}[Convolution is the bar composition;
\ClaimStatusConjectured]
\label{prop:convolution-bar}
\index{convolution product!bar composition}
\index{bar differential!convolution identification}
The convolution product on the Borel--Moore homology
$H_*^{\BM}(\Steinb_b)$ reproduces the bar differential $d_{\barB}$
and coproduct $\Delta$ of the bar complex $\barB(\Bbound)$:
\begin{enumerate}[label=\textup{(\roman*)},leftmargin=2em]
\item \textbf{Bar differential from convolution.}
  The convolution product is the composition of correspondences:
  pull back from $\Steinb_b = \Steinb_b^{(1)}$ to the triple
  product
  \[
    \Steinb_b^{(2)}
    \;=\;
    \cL_b \times_{\cM_{\mathrm{vac}}} \cL_b
    \times_{\cM_{\mathrm{vac}}} \cL_b,
  \]
  then push forward to $\Steinb_b$ via the middle projection
  $p_{13} \colon \Steinb_b^{(2)} \to \Steinb_b$ that composes
  the two outer correspondences.  Under the identification
  \eqref{eq:steinberg-factored}, the pullback lands on
  $\FM_2(\C) \times \Conf_2^{<}(\R)$, and the pushforward
  $p_{13,*}$ collapses two adjacent holomorphic points by the
  OPE residue on $\FM_2(\C)$.  This is precisely the bar
  differential $d_{\barB}$, which in degree $k$ acts by
  collapsing adjacent pairs via the chiral product.

\item \textbf{Coproduct from the diagonal.}
  The coproduct $\Delta$ comes from the reverse direction: pulling
  back from $\Steinb_b$ to $\Steinb_b^{(2)}$ via the diagonal
  map
  \[
    \delta \colon \Steinb_b \;\longrightarrow\; \Steinb_b^{(2)},
    \qquad
    (\ell_0, \ell_1)
    \;\longmapsto\;
    (\ell_0, \ell_0, \ell_1).
  \]
  Under \eqref{eq:steinberg-factored}, $\delta^*$ acts on
  $\Conf_k^{<}(\R)$ by ordered deconcatenation---splitting a
  sequence $(t_1 < \cdots < t_k)$ into
  $(t_1 < \cdots < t_j)$ and $(t_{j+1} < \cdots < t_k)$ for
  each $1 \le j \le k-1$---while leaving the $\FM_k(\C)$-factor
  unchanged.  This is the bar coproduct $\Delta$ of Vol~I.
\end{enumerate}
\end{proposition}

\begin{proof}
(i)\enspace The convolution product on $H_*^{\BM}(\Steinb_b)$ is
defined by the standard Chriss--Ginzburg recipe: given classes
$[\alpha], [\beta] \in H_*^{\BM}(\Steinb_b)$, form the external
product $[\alpha] \boxtimes [\beta] \in
H_*^{\BM}(\Steinb_b \times \Steinb_b)$, pull back to
$H_*^{\BM}(\Steinb_b^{(2)})$ via the fiber product inclusion, and
push forward via $p_{13}$.  Under the factored description
\eqref{eq:steinberg-factored} with $k = 1$, the fiber product
inclusion enforces matching of the intermediate boundary condition
at the shared $\cM_{\mathrm{vac}}$-factor, while the pushforward
$p_{13,*}$ integrates out the intermediate point.  On the
$\FM$-factor, this integration is the OPE residue (the fundamental
operation of the chiral algebra); on the $\Conf^{<}$-factor, this
integration enforces compatibility of orderings.  The result is
exactly the formula for $d_{\barB}$ in degree $k$: the sum over
$1 \le i \le k-1$ of the maps that apply the binary chiral product
$m_2$ to the $i$-th and $(i+1)$-th tensor factors.

(ii)\enspace The diagonal $\delta$ is the unit of the
correspondence calculus.  Its pullback $\delta^*$ on Borel--Moore
homology is computed via the excess intersection formula (trivial
excess bundle in the present case, since the fiber products are
transverse by the $(-2)$-shifted symplectic structure of
Proposition~\ref{prop:minus-two-shifted-general}).  Under
\eqref{eq:steinberg-factored}, $\delta^*$ acts as the identity on
$\FM_k(\C)$ (no new holomorphic points are created) and as the
Alexander--Whitney/deconcatenation map on $\Conf_k^{<}(\R)$ (the
ordered set is split at each cut point).  This is the bar coproduct
$\Delta \colon \barB_k \to \bigoplus_{j=0}^{k} \barB_j \otimes
\barB_{k-j}$, which acts by ordered deconcatenation on the
topological direction.
\end{proof}

\begin{proposition}[Swiss-cheese directionality from correspondence
geometry; \ClaimStatusConjectured]
\label{prop:directionality-correspondence}
\index{Swiss-cheese!directionality!geometric origin}
The Swiss-cheese directionality constraint---that open operations
do not produce closed outputs---is a geometric consequence of the
fiber product structure of $\Steinb_b^{(k)}$.
\end{proposition}

\begin{proof}
The coproduct $\Delta$ acts by deconcatenation on
$\Conf_k^{<}(\R)$, the topological factor of
$\Steinb_b^{(k)}$.  Crucially, $\Delta$ does not create new
points in $\FM_k(\C)$: the pullback $\delta^*$ is the identity on
the holomorphic factor.  This means the topological operation
(coproduct/open sector) cannot produce new OPE poles in the
holomorphic direction (closed sector).

Conversely, the bar differential $d_{\barB}$ collapses points in
$\FM_k(\C)$ by the OPE but preserves the cardinality and ordering
of the $\Conf^{<}(\R)$-factor (adjacent holomorphic points merge,
but their topological position is inherited by the merged point).

The fiber product over $\cM_{\mathrm{vac}}$ is the mechanism
enforcing this decoupling: it constrains the $\C$-direction
(requiring bulk field matching at each interaction point) while
the ordering on $\R$ constrains the topological direction
independently.  The two constraints commute because $\Steinb_b^{(k)}
\simeq \FM_k(\C) \times \Conf_k^{<}(\R)$ is a \emph{product},
not a fiber bundle.  The absence of open-to-closed maps is therefore
not an axiom but a theorem of the correspondence geometry: the
factored structure of $\Steinb_b^{(k)}$ forbids mixed-direction
operations that would violate the product decomposition.
\end{proof}

\begin{theorem}[Geometric Steinberg convolution identification]
\label{thm:steinberg-convolution}
\ClaimStatusConjectured
\index{Steinberg!convolution identification|textbf}
\index{Swiss-cheese convolution algebra!Steinberg identification}
There is an isomorphism of $L_\infty$-algebras
\begin{equation}\label{eq:steinberg-sc-iso}
H_*^{\BM}\bigl(\Steinb_b\bigr)
\;\cong\;
\gSC_T
\end{equation}
where $\gSC_T$ is the Swiss-cheese convolution $L_\infty$-algebra
of Definition~\textup{\ref{def:sc-convolution}}.  Under this
isomorphism:
\begin{enumerate}[label=\textup{(\roman*)},leftmargin=2em]
\item the convolution product corresponds to the bar differential
  $d_{\barB}$ \textup{(}Proposition~\textup{\ref{prop:convolution-bar}(i)}\textup{)};
\item the diagonal coproduct corresponds to the bar coproduct
  $\Delta$ \textup{(}Proposition~\textup{\ref{prop:convolution-bar}(ii)}\textup{)};
\item the Swiss-cheese directionality is enforced by the product
  decomposition of the iterated Steinberg
  \textup{(}Proposition~\textup{\ref{prop:directionality-correspondence}}\textup{)};
\item the $L_\infty$-brackets on $\gSC_T$ correspond to the higher
  convolution operations on $H_*^{\BM}(\Steinb_b^{(\bullet)})$
  coming from the full tower of iterated fiber products.
\end{enumerate}
\end{theorem}

\begin{proof}[Evidence and status]
Ingredients (a) and (b) of Conjecture~\ref{conj:geometric-steinberg}
are established in Proposition~\ref{prop:minus-two-shifted-general}
(the $(-2)$-shifted symplectic structure) and the Lagrangian
analysis of the boundary locus.  Ingredient~(c)---that convolution
on the Steinberg reproduces the Swiss-cheese algebra---is the
content of
Propositions~\ref{prop:convolution-bar}
and~\ref{prop:directionality-correspondence} above.

What remains open is the full $L_\infty$-quasi-isomorphism at the
cochain level, rather than at the level of homology.  The classical
Chriss--Ginzburg theorem identifies $H_*^{\BM}$ of the Steinberg
with the Hecke algebra as an \emph{associative} algebra; the
$L_\infty$-enrichment requires working with chains
$C_*^{\BM}(\Steinb_b^{(\bullet)})$ before passing to homology,
and establishing a homotopy transfer that is compatible with the
two-coloured operad structure.  The three worked examples
(free~$\beta\gamma$, Landau--Ginzburg, Chern--Simons) provide
evidence that this transfer exists, but the general statement
requires a derived-geometric extension of the Chriss--Ginzburg
formalism to $(-2)$-shifted symplectic stacks that is not yet
available in the literature.
\end{proof}


\section{Issues fixed, conceptual simplifications, and the platonic shape of the theory}

Six distinctions are essential for the correctness of the local theory.

\subsection*{Issue 1: identifying the wrong third vertex}
The third vertex of the triangle is not $A^!$ itself.  It is the derived center of the line category, equivalently Hochschild cochains of $A^!$.  This is fixed globally in \eqref{eq:global-corrected-triangle} and locally in Theorem~\ref{thm:local-bulk-line}.

\subsection*{Issue 2: confusing absolute and boundary-adapted line algebras}
The exact local line algebras $K_{F,p}$ and $K_\kappa$ are boundary-adapted and pointed.  They are local charts on the line sector, not the global absolute algebra $A^!$.

\subsection*{Issue 3: dependence on a chosen complement}
The intrinsic boundary algebra $B^\intr_{L,W}$ removes this at the root.  The boundary-linear exact sector then appears as the strict part of the intrinsic theory.

\subsection*{Issue 4: using homological perturbation where exact elimination is available}
Boundary Kuranishi elimination is an exact algebraic theorem, not merely a perturbative transfer.  It isolates the singular residue sharply and canonically.

\subsection*{Issue 5: hiding universal trivial blocks}
Contractible tangent-normal blocks and hyperbolic Clifford blocks are now separated cleanly.  The former are contractible; the latter are Morita-trivial after spinor stabilization.  This clarifies what should and should not count as genuine local data.

\subsection*{Issue 6: failing to distinguish pointed formal type from stable category}
The one-step Jacobi coalgebra classifies the pointed local formal germ.  The stabilization theorems explain what is invisible in the stable open sector.  The line Mather--Yau theorem classifies the reduced microlocal formal type.

\subsection*{The resulting platonic shape}

Once these corrections are imposed, the theory takes the following purified form:
\[
\text{formal boundary germ}
\;\Longrightarrow\;
B^\intr_{L,W}
\;\Longrightarrow\;
\begin{cases}
\text{exact local chart if $W$ is boundary-linear},\\
\text{full boundary $A_\infty$ algebra otherwise},
\end{cases}
\]
\[
(F,p)
\;\Longrightarrow\;
\cJ_{F,p}
\;\Longrightarrow\;
K^{\ex}_{F,p}
\;\Longrightarrow\;
K_{\kappa_F}
\;\Longrightarrow\;
\HH^\bullet(K_{\kappa_F})\cong \text{local bulk}.
\]
That is the fortified microlocal backbone from which the global theorem must be attacked.


\section{Conclusion}


On the global side, the corrected theorem-level target is clear:
\[
\Abulk
\;\simeq\;
\Zder(\Bbound)
\;\simeq\;
\Zder(\cC_{\mathrm{line}})
\;\simeq\;
\HH^\bullet_{\mathrm{ch}}(A^!),
\qquad
\cC_{\mathrm{line}}\simeq A^!\text{-mod}.
\]
On the microlocal side, the theory is no longer a network of analogies.  It is a concrete theorem package:
\begin{itemize}[leftmargin=1.8em]
\item an intrinsic boundary $A_\infty$ algebra for a formal boundary germ;
\item an exact strict dg model in the boundary-linear sector;
\item an exact pointed line algebra presented as the cobar of the one-step Jacobi coalgebra;
\item an exact Kuranishi reduction to a minimal line algebra;
\item recovery of the local bulk algebra as the derived center of that minimal line algebra;
\item and classification of the reduced microlocal formal germ by the polarized minimal line algebra.
\end{itemize}

The bulk/boundary/line Koszul triangle is therefore no longer
merely a slogan.  Its local algebraic heart has been built.  The
global structure now admits a single theorem-level formulation.

\section{The Swiss-cheese convolution algebra and the 3d universal
MC element}\label{sec:sc-convolution-mc}

\providecommand{\ModEnv}{\operatorname{ModEnv}}
\providecommand{\Tr}{\operatorname{Tr}}
\providecommand{\ohat}{\widehat{\otimes}}
\providecommand{\Mbar}{\overline{\mathcal M}}
\providecommand{\Rring}{R^{\bullet}}
\providecommand{\gSC}{\mathfrak{g}^{\mathrm{SC}}}
\providecommand{\SCop}{\mathsf{SC}^{\mathrm{ch,top}}}
\providecommand{\ov}[1]{\overline{#1}}
\providecommand{\M}{\mathcal{M}}

Volume~I constructed the modular convolution $L_\infty$-algebra
$\gAmod$ from two inputs: a chain cooperad
$C_*(\Mbar_{g,n})$ and the chiral endomorphism operad
$\End_\cA$.  The universal MC element
$\Theta_\cA \in \mc(\gAmod)$ controlled the
genus tower, and its projections recovered Theorems~A--D+H\@.

Volume~II replaces the source cooperad.  The chiral moduli spaces
$\FM_n(\C)$ and the topological intervals $E_1(m)$ assemble into
the two-coloured Swiss-cheese cooperad
$\mathbf{B}(\SCop)$---the bar construction of the homotopy-Koszul
operad of Theorem~\ref{thm:homotopy-Koszul}.
The endomorphism operad is now two-coloured:
closed operations act on the boundary algebra $\Bbound$,
open operations act on the line category $\cC_{\mathrm{line}}$,
and mixed operations encode the bulk-to-boundary map.

\begin{definition}[Swiss-cheese convolution $L_\infty$-algebra]
\label{def:sc-convolution}
\index{Swiss-cheese convolution algebra|textbf}
\index{convolution algebra!Swiss-cheese|textbf}
Let $T = (\Bbound, \cC_{\mathrm{line}})$ be an $\SCop$-algebra
on $\C \times \R_{\ge 0}$.  Define the \emph{Swiss-cheese
convolution $L_\infty$-algebra} by
\begin{equation}\label{eq:gSC-def}
\gSC_T
\;:=\;
\Hom_{\Sigma}\!\bigl(
  \mathbf{B}(\SCop),\;
  \End_{(\Bbound,\,\cC_{\mathrm{line}})}
\bigr),
\end{equation}
the two-coloured convolution of the bar cooperad with the
two-coloured endomorphism operad.  The $L_\infty$-brackets are
indexed by two-coloured stable trees: at each closed vertex the
chiral operation $\Phi_v^{\mathrm{ch}}$ is inserted, at each open
vertex the topological operation $\Phi_v^{\mathrm{top}}$ is
inserted, and at each mixed vertex the bulk-to-boundary operation
is inserted.
\end{definition}

The key structural fact: the two-coloured bar cooperad
$\mathbf{B}(\SCop)$ is \emph{not} a product of a closed and an
open cooperad.  It is a single two-coloured object whose mixed
operations encode the exact correlation between holomorphic
collisions on $\C$ and topological orderings on $\R$.  The
Swiss-cheese convolution algebra $\gSC_T$ therefore carries
strictly more information than the pair $(\gAmod[\Bbound],
\End(\cC_{\mathrm{line}}))$: it remembers the mixed
bulk-boundary-line interactions.

\begin{theorem}[3d universal MC element; \ClaimStatusProvedHere]
\label{thm:3d-universal-mc}
\index{universal MC element!3d|textbf}
\index{Maurer--Cartan element!Swiss-cheese|textbf}
Let $T = (\Bbound, \cC_{\mathrm{line}})$ be an $\SCop$-algebra
satisfying~\textup{(H1)--(H4)}.  There exists a canonical
Maurer--Cartan element
\begin{equation}\label{eq:alpha-T}
\alpha_T
\;\in\;
\mc\bigl(\gSC_T\bigr)
\end{equation}
whose two-coloured graph expansion is
\[
\alpha_T
\;=\;
\sum_{\Gamma \in \mathsf{Gr}^{\mathrm{SC}}}
\frac{1}{|\operatorname{Aut}(\Gamma)|}
\;W_\Gamma\;\otimes\;\Phi_\Gamma^T,
\]
where $\mathsf{Gr}^{\mathrm{SC}}$ ranges over connected
two-coloured stable trees with closed vertices labelled by
$\FM_k(\C)$-chains, open vertices labelled by $E_1(m)$-chains,
and mixed vertices labelled by $\FM_k(\C) \times E_1(m)$-chains.
The MC equation
$D\alpha_T + \tfrac{1}{2}[\alpha_T, \alpha_T] = 0$
holds because $\SCop$ is homotopy-Koszul
\textup{(}Theorem~\textup{\ref{thm:homotopy-Koszul}}\textup{)},
hence $\mathbf{B}(\SCop)$ has $\partial^2 = 0$.
\end{theorem}

\begin{proof}
By Theorem~\ref{thm:homotopy-Koszul}, $\SCop$ is homotopy-Koszul:
$\Omega\mathbf{B}(\SCop) \xrightarrow{\sim} \SCop$.  The universal
twisting morphism
$\tau\colon \mathbf{B}(\SCop) \to \SCop$ therefore satisfies the
MC equation in the two-coloured convolution algebra.  Postcomposing
$\tau$ with the structure map $\SCop \to \End_{(\Bbound,
\cC_{\mathrm{line}})}$ gives $\alpha_T$, and the MC equation is
preserved under this postcomposition.  The graph expansion follows
from the definition of the two-coloured bar construction: each
two-coloured tree $\Gamma$ contributes a term $W_\Gamma \otimes
\Phi_\Gamma^T$, where $W_\Gamma$ is the fundamental chain of the
corresponding boundary stratum of $\FM_k(\C) \times E_1(m)$ and
$\Phi_\Gamma^T$ is the composition of structure operations along
the tree.
\end{proof}

The element $\alpha_T$ is the 3d analogue of the Vol~I element
$\Theta_\cA$.  Its projections recover every structure of Vol~II:

\begin{proposition}[Projections of $\alpha_T$;
\ClaimStatusProvedHere]
\label{prop:alpha-projections}
\index{universal MC element!3d projections}
The MC element $\alpha_T$ of
Theorem~\textup{\ref{thm:3d-universal-mc}} determines:
\begin{enumerate}[label=\textup{(\roman*)}]
\item \emph{Closed face.}  The restriction of $\alpha_T$ to
  closed vertices recovers the boundary chiral algebra~$\Bbound$
  with its $\Ainf$-structure
  \textup{(}Definition~\textup{\ref{def:ainfty_chiral}}\textup{)}.
\item \emph{Open face.}  The restriction to open vertices
  recovers the line category $\cC_{\mathrm{line}} \simeq
  A^!\text{-mod}$ with its $\Ainf$-module structure
  \textup{(}\S\textup{\ref{sec:line_operators}}\textup{)}.
\item \emph{Mixed face: spectral $R$-matrix.}  The arity-$(2,0)$
  component of $\alpha_T$---two closed vertices, no open
  vertices---restricted to the ordered collision stratum of
  $\FM_2(\C)$ is the spectral braiding $R(z)$ of
  Definition~\textup{\ref{def:spectral-braiding}}.  Its MC
  equation is the Yang--Baxter equation
  \textup{(}Theorem~\textup{\ref{thm:YBE}}\textup{)}.
\item \emph{Mixed face: bulk-to-boundary.}  The arity-$(1,1)$
  component---one closed vertex, one open vertex---is the
  bulk-to-boundary map $\Abulk \to \Zder(\Bbound)$.  Its MC
  equation is the factorization constraint
  \textup{(}Theorem~\textup{\ref{thm:boundary-linear-bulk-boundary}}\textup{)}.
\item \emph{Genus face.}  At genus~$g \ge 1$, the restriction
  of $\alpha_T$ to the closed sector recovers the modular
  characteristic $\kappa(\Bbound) \cdot \omega_g$ and the
  shadow Postnikov tower $\Theta_{\Bbound}^{\le r}$ of Vol~I.
\item \emph{PVA descent.}  Passing to cohomology on the closed
  face produces the Poisson vertex algebra bracket
  $\{-_\lambda-\}$ on $H^\bullet(\Bbound)$
  \textup{(}Theorem~\textup{\ref{thm:cohomology_PVA}}\textup{)}.
\end{enumerate}
\end{proposition}

\begin{proof}
Each projection is the restriction of the two-coloured graph sum
to graphs of a specific colour type.

(i): Graphs with all vertices closed are the chiral trees
$\FM_n(\C)$; their operations are the $\Ainf$-chiral operations
$m_k$ by the axioms of the $\SCop$-algebra structure.

(ii): Graphs with all vertices open are the $E_1$-trees; their
operations define the $\Ainf$-module structure on lines by the
homotopy-Koszul property.

(iii): The $(2,0)$-graph is a single edge connecting two closed
vertices at positions $z_1, z_2 \in \FM_2(\C)$.  The structure
operation is $\mathsf{Comp}_{z_1,z_2}$, which is the spectral
braiding $R(z_{12})$ by Definition~\ref{def:spectral-braiding}.
The MC equation $D\alpha + \tfrac{1}{2}[\alpha,\alpha] = 0$
restricted to arity $(3,0)$ gives the YBE by the same
boundary-face argument as Theorem~\ref{thm:YBE}.

(iv): The $(1,1)$-graph is a single mixed edge.  The MC equation
at this colour type is the factorization-module consistency, which
by Theorem~\ref{thm:boundary-linear-bulk-boundary} gives the
bulk-as-derived-center identification.

(v): Closing the genus face by including self-loop graphs in the
closed sector reproduces the modular bar construction of Vol~I;
the MC element restricted to closed self-loop graphs at genus~$g$
is $\Theta_\cA|_{g}$ by the same argument as
Theorem~\ref*{thm:mc2-bar-intrinsic} (bar-intrinsic construction).

(vi): Descent to cohomology on the closed face is the content of
Theorem~\ref{thm:cohomology_PVA}: the $\Ainf$-operations $m_k$
descend to the PVA bracket $\{-_\lambda-\}$ by the repaired
descent argument of \S\ref{sec:PVA_descent}.
\end{proof}

\begin{definition}[3d modular homotopy type]
\label{def:3d-modular-homotopy-type}
\index{modular homotopy type!3d|textbf}
The \emph{3d modular homotopy type} of a holomorphic-topological
theory~$T$ is the formal moduli problem
\[
\cM^{\mathrm{SC}}_T(R)
\;:=\;
\mc\bigl(\gSC_T \ohat R\bigr)\big/{\sim}
\]
for local Artin dg algebras~$R$.  Its basepoint is $\alpha_T$.
\end{definition}

\begin{corollary}[Restriction to Vol~I; \ClaimStatusProvedHere]
\label{cor:3d-to-vol1-restriction}
The restriction of $\alpha_T$ to closed vertices with
genus expansion recovers the Vol~I universal MC element:
\[
\alpha_T\big|_{\text{closed, all genera}}
\;=\;
\Theta_{\Bbound}.
\]
The 3d modular homotopy type $\cM^{\mathrm{SC}}_T$ refines the
Vol~I modular homotopy type $\cM^{\mathrm{mod}}_{\Bbound}$:
the forgetful map $\gSC_T \to \gAmod[\Bbound]$ that discards
open and mixed vertices induces a morphism of formal moduli
problems $\cM^{\mathrm{SC}}_T \to \cM^{\mathrm{mod}}_{\Bbound}$
whose fibre over $\Theta_{\Bbound}$ is the space of line-category
deformations compatible with the boundary algebra.
\end{corollary}


\section{Kodaira--Spencer holography and modular defect observables}
\label{sec:ks-holography}

The preceding sections built the local microlocal engine.
We now connect that engine to a concrete large-$N$ holographic
package---Kodaira--Spencer gravity coupled to holomorphic
Chern--Simons---and extract modular descendants and defect indices
that organize the higher-genus output.

\subsection{The Kodaira--Spencer universal defect algebra}
\label{subsec:ks-defect-algebra}

\begin{example}[Large-$N$ boundary algebra and universal defect generators;
  \ClaimStatusProvedElsewhere]
\label{def:ks-large-n-boundary}
In the twisted-holography package of \cite{Zeng23} and~\cite{KZ25},
the large-$N$ open-side chiral algebra is built from the following
fields:
\begin{itemize}[leftmargin=1.8em]
\item $(b,c)$ a ghost system valued in $\mathfrak{gl}_N$;
\item $(Z_1,Z_2)$ an adjoint $\beta\gamma$ system;
\item $(I,J)$ a fundamental $\beta\gamma$ system;
\end{itemize}
with BRST operator
\[
Q_{\mathrm{BRST}}
\;=\;
\oint\!\Bigl(\Tr(c\,Z_1 Z_2)
  \;+\;\Tr(I\,c\,J)
  \;+\;\tfrac{1}{2}\,\Tr(b\,c^2)\Bigr).
\]
Its large-$N$ BRST cohomology contains towers
$A^{(n)},\;B^{(n)},\;C^{(n)},\;D^{(n)},\;E_t^{(n)}$,
where the first four match Kodaira--Spencer gravity and the
$E$-tower matches holomorphic Chern--Simons.
\end{example}

\begin{definition}[Kodaira--Spencer universal defect algebra]
\label{def:ks-universal-defect}
Define the \emph{Kodaira--Spencer universal defect algebra} to be
\[
\mathcal{U}_{\mathrm{KS}}
\;:=\;
V_{\partial,\mathrm{KS}}(B_{\mathrm{univ}}),
\]
with subalgebras:
\begin{itemize}[leftmargin=1.8em]
\item \emph{Gravitational sector:}
  $\mathcal{U}_{\mathrm{KS}}^{\mathrm{cl}}
  :=\langle A^{(n)},\,B^{(n)},\,C^{(n)},\,D^{(n)}\rangle$;
\item \emph{Gauge sector:}
  $\mathcal{U}_{\mathrm{KS}}^{\mathrm{gauge}}
  :=\langle E_t^{(n)}\rangle$.
\end{itemize}
The modular state spaces are defined by factorization integration:
\[
\mathcal{H}_{g,n}^{\mathrm{KS}}
\;:=\;
\int_{\Sigma_{g,n}}\!\mathcal{U}_{\mathrm{KS}},
\qquad
\mathcal{H}_{g,n}^{\mathrm{cl,KS}}
\;:=\;
\int_{\Sigma_{g,n}}\!\mathcal{U}_{\mathrm{KS}}^{\mathrm{cl}}.
\]
\end{definition}

\begin{conjecture}[Open--closed modular splitting for KS holography;
  \ClaimStatusConjectured]
\label{conj:ks-open-closed-splitting}
There is a functorial completed factorization
\[
\mathcal{H}_{g,n}^{\mathrm{KS}}
\;\simeq\;
\mathcal{H}_{g,n}^{\mathrm{cl,KS}}
\;\ohat\;
\mathcal{H}_{g,n}^{\mathrm{gauge,KS}}
\]
compatible with the holographic bulk--boundary map and with stable-curve
gluing
$\mu_{g,n}^{\mathrm{KS}}$.
Precisely: there is an equivalence of dg categories
$\cC_{\mathrm{line}}^{\mathrm{KS}} \simeq
\cC_{\mathrm{line}}^{\mathrm{cl,KS}} \boxtimes
\cC_{\mathrm{line}}^{\mathrm{gauge,KS}}$
(external tensor product), and the factorization-homology functor preserves
this splitting:
$\int_{\Sigma_{g,n}} (-) \simeq
\int_{\Sigma_{g,n}} (-)|_{\mathrm{cl}} \;\ohat\;
\int_{\Sigma_{g,n}} (-)|_{\mathrm{gauge}}$.
This gives a mathematically sharp candidate for separating pure
Kodaira--Spencer gravity from holomorphic Chern--Simons at all genera
within a single universal defect algebra.
\end{conjecture}


\subsection{Descendant potentials and defect indices}
\label{subsec:descendant-potentials}

\begin{definition}[Holographic tautological action]
\label{def:descendant-taut-action}
Suppose $\ModEnv(\cT;B)$ exists.  The \emph{universal correlator class} is
\[
\mathcal{U}_{g,n}^{\cT;B}
\;\in\;
H^{\bullet}\!\bigl(\Mbar_{g,n};\;\End(\mathcal{H}_{g,n}^{\cT;B})\bigr).
\]
The \emph{holographic tautological action} is the map
\[
\Theta_{g,n}^{\cT;B}
\;:\;
\Rring(\Mbar_{g,n})
\;\longrightarrow\;
\End(\mathcal{H}_{g,n}^{\cT;B}),
\qquad
\alpha\;\longmapsto\;
\int\alpha\cup\mathcal{U}_{g,n}^{\cT;B},
\]
where $\Rring(\Mbar_{g,n})$ denotes the tautological ring.
\end{definition}

\begin{definition}[Modular holographic descendant potential]
\label{def:descendant-potential}
The \emph{modular holographic descendant potential} is
\[
\mathcal{F}_{\cT;B}(\hbar,t,s)
\;:=\;
\sum_{g,n}\;\sum_{a_1,\dots,a_n}\;\sum_{k_1,\dots,k_n\ge 0}
\frac{\hbar^{2g-2+n}}{n!}\;
\bigl\langle
  \textstyle\prod_{i=1}^{n}\tau_{k_i}(\cO_{a_i})
\bigr\rangle_{g,n}^{\cT;B}\;
\prod_{i=1}^{n} t_{a_i}\,s_{k_i},
\]
where the descendants
$\tau_{k}(\cO_{a})$ are defined through the tautological action
$\Theta_{g,n}^{\cT;B}$.
\end{definition}

\begin{definition}[Modular defect index]
\label{def:defect-index}
When an $L_0$-grading is defined on
$\mathcal{H}_{g,n}^{\cT;B}$, the \emph{modular defect index} is
\[
\mathcal{I}_{\cT;B}^{\mathrm{def}}(q,\hbar)
\;:=\;
\sum_{g,n\ge 0}\;
\hbar^{2g-2+n}\;
\Tr_{\mathcal{H}_{g,n}^{\cT;B}}\!\bigl((-1)^F\,q^{L_0}\bigr).
\]
\end{definition}

\begin{conjecture}[Tautological recursion for holographic descendants;
  \ClaimStatusConjectured]
\label{conj:taut-recursion}
The potential $\mathcal{F}_{\cT;B}$ satisfies all tautological
relations imposed by the image of
$\Rring(\Mbar_{g,n})$ under
$\Theta_{g,n}^{\cT;B}$.  When the tautological action $\Theta_{g,n}^{\cT;B}$ factors through a
semisimple cohomological field theory (CohFT), the descendant potential
$\mathcal{F}_{\cT;B}$ satisfies the Givental--Teleman classification and is
uniquely determined by genus-$0$ data via $R$-matrix action.  In
non-semisimple cases (e.g., logarithmic theories), the potential is expected
to satisfy the Eynard--Orantin topological recursion with spectral curve
determined by the genus-$0$ two-point function $r_{T,0}(z)$.
\end{conjecture}


\appendix

\section{Notation guide}

\begin{center}
\renewcommand{\arraystretch}{1.2}
\begin{longtable}{|p{0.28\textwidth}|p{0.6\textwidth}|}
\hline
$L\subset V$ & boundary subspace inside ambient linear space \\
\hline
$Q=V/L$ & intrinsic normal quotient \\
\hline
$W$ & formal superpotential vanishing on $L$ \\
\hline
$B^\intr_{L,W}$ & intrinsic boundary $A_\infty$ algebra \\
\hline
$F:L\to N^\vee$ & formal map defining the boundary-linear potential $W(x,y)=\langle y,F(x)\rangle$ \\
\hline
$B_{L,W}$ & strict boundary dg algebra in the boundary-linear sector \\
\hline
$Z_F^{\der}$ & derived zero locus of $F$ \\
\hline
$\dCrit(W)$ & derived critical locus of the bulk potential \\
\hline
$\cJ_{F,p}$ & one-step Jacobi coalgebra of the pointed germ $(F,p)$ \\
\hline
$K^{\ex}_{F,p}$ & exact pointed line algebra \\
\hline
$\kappa$ & reduced Kuranishi map after eliminating linearly massive directions \\
\hline
$K_\kappa$ & minimal pointed line algebra \\
\hline
$A^!$ & absolute Koszul-dual algebra conjecturally controlling all lines in the global theory \\
\hline
$\Zder$ & derived center \\
\hline
$\HH^\bullet$ & Hochschild cochains/cohomology \\
\hline
\end{longtable}
\end{center}

% Bibliography entries consolidated into main.tex.


%================================================================
% THE GENUS-ZERO HOLOGRAPHIC SPINE AND AMBIENT COMPLEMENTARITY
%================================================================

\section{The genus-zero holographic spine and ambient complementarity}
\label{sec:genus-zero-holographic-spine}

The bulk-boundary-line triangle of the preceding sections
is a genus-zero snapshot.  We now show that it is the \emph{local face}
of a single ambient Maurer--Cartan element, and that the correct
object is not a modular envelope together with corrections, but one
ambient Hamiltonian deformation problem whose five-piece differential
encodes sewing, collision, genus-raising, and Koszul twisting
simultaneously.  Costello's $M2$-brane example is the first
complete realization.

\subsection{Established input: the genus-zero package}
\label{subsec:established-genus-zero-package}

The papers most directly governing the construction are the
following.

\begin{theorem}[Genus-zero holographic spine]
\label{thm:genus-zero-holographic-spine}
\ClaimStatusProvedElsewhere
Let $\mathcal{T}$ be a perturbative holomorphic-topological theory
with a rich boundary or defect sector.  Then:
\begin{enumerate}[label=\textup{(\arabic*)}]
\item Bulk local operators form a commutative chiral algebra with
shifted Poisson bracket and higher stress tensor; boundary local
operators form a generally noncommutative chiral algebra with
bulk-boundary map and higher $A_\infty$ operations
\textup{(Costello--Dimofte--Gaiotto)}.
\item For KK reductions of six-dimensional holomorphic theories,
the reduced theory is analysed in BV form, effective interactions
are produced at all classical orders by homotopy transfer, and
boundary conditions identify the boundary chiral algebra with the
universal defect chiral algebra \textup{(Zeng)}.
\item Higher operations are computed by regularized Feynman diagrams
and satisfy quadratic consistency relations interpretable as higher
Wess--Zumino conditions \textup{(Gaiotto--Kulp--Wu)}.
\item The category of perturbative lines is equivalent to modules
for a Koszul-dual $A_\infty$-algebra $A^!$, and the OPE of lines
is controlled by a dg-shifted Yangian with Maurer--Cartan kernel
$r(z) \in A^! \widehat\otimes A^!(\!(z^{-1})\!)$
solving an $A_\infty$-Yang--Baxter equation
\textup{(Dimofte--Niu--Py)}.
\item In twisted holography, Costello's $11$-dimensional
$\Omega$-background admits a five-dimensional noncommutative
Chern--Simons realization, and in the $M2$-brane example the
large-$K$ operator algebra is Koszul dual to the bulk algebra
and becomes a quantum double-loop algebra \textup{(Costello)}.
\end{enumerate}
\end{theorem}

\begin{remark}
This theorem should be read as the exact reason that the correct
object cannot be merely a modular envelope.  The literature already
forces four structures to coexist: cyclic deformation theory, Koszul
duality, modular sewing, and meromorphic line-operator transport.
A bare modular envelope sees only the sewing part.
\end{remark}

\begin{definition}[Genus-zero package]
\label{def:genus-zero-package}
For a holomorphic-topological theory $\mathcal{T}$ with boundary
condition $B$, the genus-zero package is the tuple
\begin{equation}
\label{eq:genus-zero-package}
G_0(\mathcal{T}; B)
:=
\bigl(
V_{\mathcal{T}},\,
\{\!\{-,-\}\!\},\,
V_{\partial,\mathcal{T}}(B),\,
\beta_B,\,
A^!_{\mathcal{T}},\,
\mathcal{C}_{\mathcal{T}},\,
r_{\mathcal{T}}(z),\,
\{m_k^{\mathcal{T}}\}_{k \geq 2}
\bigr),
\end{equation}
where $V_{\mathcal{T}}$ is the bulk chiral algebra,
$V_{\partial,\mathcal{T}}(B)$ the boundary chiral algebra,
$\beta_B$ the bulk-boundary map, $A^!_{\mathcal{T}}$ a Koszul dual,
$\mathcal{C}_{\mathcal{T}} \simeq A^!_{\mathcal{T}}\text{-Mod}$
the line category, and $r_{\mathcal{T}}(z)$ the line-OPE
Maurer--Cartan kernel.
\end{definition}

\subsection{The ambient complementarity principle}
\label{subsec:ambient-complementarity-principle}

\begin{principle}[Ambient complementarity]
\label{princ:ambient-complementarity}
The true object of the theory is one ambient deformation problem in
which the $A$-side and the $A^!$-side are opposite Lagrangians.
Modular gluing, higher operations, line-operator transport, and
resonance all live inside that one ambient problem.
\end{principle}

The decisive shift is: stop thinking of complementarity as an
additive splitting and instead regard it as a \emph{polarization}
of a shifted-symplectic deformation problem.  The complementarity
potential $S_{\mathcal{A}}$
(Definition~\ref*{def:complementarity-potential} in Vol~I) is the
generating function of one Lagrangian in the cotangent chart
centered at the other.  Its Taylor jets are the shadow tensors
$H_{\mathcal{A}}, \mathfrak{C}_{\mathcal{A}},
\mathfrak{Q}_{\mathcal{A}}$ computed in
Vol~I for each standard family.

\begin{conjecture}[Modular completion of the genus-zero package]
\label{conj:modular-completion-package}
\ClaimStatusConjectured
Every package $G_0(\mathcal{T}; B)$ arising from the
holographic/twisted-holomorphic setting admits a modular completion
\[
\mu^{\mathcal{T}}_{g,n} \colon
C_\bullet(\overline{\mathcal{M}}_{g,n+1})
\otimes (A^!_{\mathcal{T}})^{\widehat\otimes n}
\to A^!_{\mathcal{T}},
\qquad 2g - 2 + n > 0,
\]
whose genus-zero truncation is $\{m_n^{\mathcal{T}}\}$, whose
two-leg projection is the modular expansion of
$r_{\mathcal{T}}(z)$, and whose generating Maurer--Cartan element
is precisely $\Theta_{\mathcal{A}}$ for the ambient
complementarity algebra attached to $A_{\partial,\mathcal{T}}$.
\end{conjecture}

\begin{construction}[From genus-zero package to holographic
modular Koszul datum; \ClaimStatusProvedHere]
\label{constr:genus-zero-to-hmkd}
\index{holographic modular Koszul datum!from genus-zero package}
The genus-zero package $G_0(\mathcal T;B)$ embeds into the
holographic modular Koszul datum
$\mathcal H(T)=(\cA,\cA^!,\cC,r(z),\Theta_\cA,
\nabla^{\mathrm{hol}})$
of Vol~I (Definition~\ref*{def:holographic-modular-koszul-datum})
by the following identifications:
\begin{center}
\renewcommand{\arraystretch}{1.15}
\begin{tabular}{ll}
\textbf{Genus-$0$ package} & \textbf{Holographic datum} \\
\hline
$V_\mathcal T$ (bulk chiral)
  & $\cA$ \\
$V_{\partial,\mathcal T}(B)$ (boundary)
  & $\cA$ restricted to $\{0\}\times X$ \\
$A^!_\mathcal T$ (Koszul dual)
  & $\cA^!$ \\
$\cC_\mathcal T$ (line category)
  & $\cC\simeq\cA^!\text{-mod}$ \\
$r_\mathcal T(z)$ (line-OPE kernel)
  & $r(z)=\operatorname{Res}^{\mathrm{coll}}_{0,2}(\Theta_\cA)$ \\
$\{m_k^\mathcal T\}_{k\ge 2}$ (higher operations)
  & $\{\ell_k^{(0)}\}_{k\ge 2}$ (genus-$0$ $L_\infty$ brackets) \\
--- & $\Theta_\cA$ (modular MC element) \\
--- & $\nabla^{\mathrm{hol}}_{g,n}$ (shadow connection)
\end{tabular}
\end{center}
The last two entries are the modular extension: they have no
genus-$0$ counterpart and arise from the bar-intrinsic
construction (Vol~I,
Theorem~\ref*{thm:mc2-bar-intrinsic}).

The modular extension adds:
\begin{enumerate}[label=\textup{(\roman*)}]
\item \emph{Genus-$1$ curvature.}
  $d_{\mathrm{fib}}^2=\kappa(\cA)\cdot\omega_g$: the
  one-loop correction to the bar differential.  The genus-$0$
  package sees $d^2=0$; the modular extension sees
  $d^2=\kappa\cdot\omega_g$.
\item \emph{The genus tower.}
  $F_g(\cA)=\kappa\cdot B_{2g}/(2g(2g-2))$: the scalar
  part of the genus-$g$ free energy.
\item \emph{The shadow Postnikov tower.}
  $\Theta_\cA^{\le 2}\to\Theta_\cA^{\le 3}\to\cdots$:
  finite-order projections carrying nonlinear
  information invisible at genus~$0$.
\item \emph{The shadow connection.}
  $\nabla^{\mathrm{hol}}_{g,n}=d-\mathrm{Sh}_{g,n}(\Theta_\cA)$:
  the flat connection on derived conformal blocks.  At
  $g=0$ this is the system of differential equations
  determining correlators (KZ for affine, BPZ for Virasoro);
  at $g\ge 1$ it acquires curvature $\kappa\cdot\omega_g$.
\end{enumerate}
\end{construction}

\begin{computation}[Modular extension of the Heisenberg
genus-zero package; \ClaimStatusProvedHere]
\label{comp:vol2-heisenberg-modular-extension}
\index{Heisenberg algebra!modular extension}
The Heisenberg genus-$0$ package has
$r(z)=k/z$, $m_k=0$ for $k\ge 3$, and $\cC=\mathrm{Vect}$.
The modular extension:
$\Theta_{\cH_k}=k\cdot\eta\otimes\Lambda$,
$\kappa=k$,
$F_g=k\cdot B_{2g}/(2g(2g-2))$,
$\nabla^{\mathrm{hol}}_{0,n}
=d-k\sum_{i<j}q_iq_j\,d\log(z_i-z_j)$
(flat at all genera).
Shadow tower terminates at arity~$2$; all nonlinear
shadows vanish.
\end{computation}

\begin{computation}[Modular extension of the affine
genus-zero package; \ClaimStatusProvedHere]
\label{comp:vol2-affine-modular-extension}
\index{Kac--Moody!modular extension}
For $\widehat{\mathfrak{sl}}_2$ at level~$k$: the genus-$0$
package has $r(z)=\Omega/z$, $m_3\neq 0$ (cubic from Jacobi),
$m_k=0$ for $k\ge 4$ (strict formality).
The modular extension adds:
$\kappa=3k/(k+2)$,
$\nabla^{\mathrm{hol}}_{0,n}$ = KZ connection
(derived in Vol~I,
Computation~\ref*{comp:kz-from-graph-sum}),
$\nabla^{\mathrm{hol}}_{1,1}=d-\kappa\,\omega_1$
(genus-$1$ curvature from Hodge class),
spectral discriminant
$\Delta(x)=(1-kx)(1-(k{+}4)x)/(1-2x)$.
Shadow depth~$3$; no planted-forest corrections.
\end{computation}

\begin{computation}[Modular extension of the Virasoro
genus-zero package; \ClaimStatusProvedHere]
\label{comp:vol2-virasoro-modular-extension}
\index{Virasoro algebra!modular extension}
For $\mathrm{Vir}_c$: the genus-$0$ package has
$r_c(z)=\partial T/z+2T/z^2+c/(2z^4)$, $m_3\neq 0$, $m_4\neq 0$
(quartic contact $10/[c(5c+22)]$).
The modular extension adds:
$\kappa=(c-26)/2$,
$\nabla^{\mathrm{hol}}_c
=d-\frac{c-26}{2}\,\omega_g
-\frac{10}{c(5c+22)}\,\delta_4
-\frac{120}{c^2(5c+22)}\,\delta_H-\cdots$,
genus-$2$ three-shell decomposition, infinite shadow depth.
At $c=26$: $\kappa=0$, shadow connection flat at leading order.
\end{computation}

\begin{computation}[Shadow connection as hypergeometric ODE:
$\widehat{\mathfrak{sl}}_2$ at $4$ points;
\ClaimStatusProvedHere]
\label{comp:shadow-connection-sl2-4pt}
\index{shadow connection!hypergeometric ODE}
\index{KZ connection!4-point hypergeometric}
Fix four $\widehat{\mathfrak{sl}}_2$-insertions at
$z_1=0$, $z_2=z$, $z_3=1$, $z_4=\infty$ with spins
$j_1=j_2=j_3=j_4=1/2$ (all fundamental).
The $4$-point conformal block $F(z)$ satisfies
$\nabla^{\mathrm{hol}}_{0,4}\,F=0$, which reduces to
\begin{equation}\label{eq:kz-hypergeometric}
\frac{dF}{dz}
\;=\;
\frac{1}{k+2}
\left(
\frac{\Omega_{12}}{z}
+\frac{\Omega_{23}}{z-1}
\right)F(z),
\end{equation}
where $\Omega_{12}$ and $\Omega_{23}$ act on
$V^{\otimes 4}\cong(\mathbb C^2)^{\otimes 4}$.

Decompose: $V\otimes V=V_0\oplus V_1$ (singlet $\oplus$
triplet).  The $4$-point block in the $s$-channel
$((V\otimes V)\otimes V)\otimes V$ decomposes as
$F(z)=\alpha\,f_0(z)+\beta\,f_1(z)$
where $f_0$ and $f_1$ are the conformal blocks in the
singlet and triplet intermediate channels.

The ODE~\eqref{eq:kz-hypergeometric} reduces to
a scalar hypergeometric equation on each channel:
\begin{align*}
f_0(z)
&\;=\;
z^{-\lambda_0}(1-z)^{-\mu_0}\,
{}_2F_1\!\bigl(a_0,b_0;c_0;\,z\bigr),\\
f_1(z)
&\;=\;
z^{-\lambda_1}(1-z)^{-\mu_1}\,
{}_2F_1\!\bigl(a_1,b_1;c_1;\,z\bigr),
\end{align*}
where $\lambda_i=\Omega^{(i)}_{12}/(k{+}2)$,
$\mu_i=\Omega^{(i)}_{23}/(k{+}2)$, and
$\Omega^{(i)}$ denotes the Casimir eigenvalue
on the $i$-th intermediate channel.

At $k=1$ ($q=e^{i\pi/3}$): $\lambda_0=1/6$,
$\lambda_1=-1/6$, and the monodromy of $F(z)$ around
$z=0$ is the $R$-matrix
$R=\operatorname{diag}(e^{-i\pi/3},\,e^{i\pi/3})$
in the singlet/triplet basis.

At $k\to\infty$ (classical limit): $\lambda\to 0$ and
the hypergeometric functions degenerate to constants
(free propagation, no scattering).
\end{computation}

\begin{computation}[Genus-$1$ partition functions
from the shadow connection; \ClaimStatusProvedHere]
\label{comp:genus1-partition-functions}
\index{partition function!genus 1}
\index{shadow connection!genus 1}
The genus-$1$ partition function $Z_1(\cA)$
is the trace of the genus-$1$ shadow operator
$\operatorname{tr}_{V_1}(q^{L_0-c/24})$ where
$V_1$ is the space of genus-$1$ conformal blocks.
At the scalar level, $F_1=\kappa\cdot B_2/2=-\kappa/12$.
\begin{center}
\renewcommand{\arraystretch}{1.2}
\begin{tabular}{llll}
\textbf{Family} & $\kappa$ & $F_1=-\kappa/12$
  & \textbf{Partition function} \\
\hline
$\cH_k$ & $k$ & $-k/12$
  & $\eta(q)^{-k}$ \\
$\widehat{\mathfrak{sl}}_2{}_k$
  & $3k/(k{+}2)$ & $-k/(4(k{+}2))$
  & $\chi_k(q)$ (level-$k$ character) \\
$\mathrm{Vir}_c$ & $(c{-}26)/2$ & $-(c{-}26)/24$
  & $q^{-(c-1)/24}\prod_n(1-q^n)^{-1}$ \\
$\beta\gamma$ & $-2$ & $1/6$
  & $\eta(q)^2$ \\
$\mathcal W_3{}_c$ & $(c{-}50)/2$ & $-(c{-}50)/24$
  & $q^{-(c-2)/24}\prod_n(1-q^n)^{-2}$
\end{tabular}
\end{center}
The $\eta(q)^{-k}$ for Heisenberg is the Dedekind eta
function to the power $-k$: one bosonic oscillator per unit
of level.  The sign of $F_1$ for $\beta\gamma$ ($F_1>0$)
reflects $\kappa<0$ (fermionic genus tower).
\end{computation}

\begin{construction}[Boundary from bulk via derived center;
\ClaimStatusProvedHere]
\label{constr:boundary-from-bulk}
\index{derived center!boundary from bulk}
Given a bulk $E_2$-algebra $B$, the boundary algebra is the
derived center
$Z(B):=R\!\operatorname{Hom}_{B^e}(B,B)$
where $B^e=B\otimes B^{\mathrm{op}}$.  For the standard families:
\begin{center}
\renewcommand{\arraystretch}{1.15}
\begin{tabular}{lll}
\textbf{Bulk $B$} & \textbf{$Z(B)$} & \textbf{Boundary} \\
\hline
$\Bbbk[[\kappa]]$ & $\cH_k$ & Heisenberg \\
$\Bbbk[[\Omega]]$ & $\widehat{\mathfrak{sl}}_2{}_k$ & affine \\
$\Bbbk[[c]]\otimes\operatorname{Sym}^*(L_n^*)$
  & $\mathrm{Vir}_c$ & Virasoro \\
$\Bbbk[[c]]\otimes\operatorname{Sym}^*(L_n^*\oplus W_m^*)$
  & $\mathcal W_3{}_c$ & $\mathcal W_3$
\end{tabular}
\end{center}
At the critical level $k=-h^\vee$: $Z(B)$ acquires the
Feigin--Frenkel center as an additional factor, and the
boundary algebra jumps.
\end{construction}

\begin{computation}[Genus-$2$ partition function: first
computation; \ClaimStatusProvedHere]
\label{comp:genus2-partition-function}
\index{partition function!genus 2}
The genus-$2$ free energy $F_2(\cA)$ decomposes via the
three-shell formula
$\Theta^{(2)}=\Theta^{(2)}_{\mathrm{loop}^2}
+\Theta^{(2)}_{\mathrm{sep}\circ\mathrm{loop}}
+\Theta^{(2)}_{\mathrm{pf}}$.

At the scalar level:
$F_2=\kappa^2\cdot B_4/(4\cdot 2)=\kappa^2/(240)$
(from $B_4=-1/30$ and the Bernoulli normalisation).

\begin{center}
\renewcommand{\arraystretch}{1.2}
\begin{tabular}{lll}
& $F_2=\kappa^2/240$ & shells \\
\hline
$\cH_k$ & $k^2/240$ & loop$^2$ only \\
$\widehat{\mathfrak{sl}}_2{}_k$
  & $9k^2/(240(k{+}2)^2)$
  & loop$^2$+sep$\circ$loop \\
$\mathrm{Vir}_c$ & $(c{-}26)^2/960$
  & all three \\
$\beta\gamma$ & $4/240=1/60$ & loop$^2$ only \\
$\mathcal W_3{}_c$ & $(c{-}50)^2/960$
  & all three
\end{tabular}
\end{center}
For $\mathrm{Vir}_{26}$: $F_2=0$ (anomaly cancellation extends
to genus~$2$).  For $\mathcal W_3{}_{50}$: $F_2=0$ likewise.
The non-scalar corrections (from the sep$\circ$loop and
planted-forest shells) add terms proportional to the cubic
shadow $\mathfrak C$ and the quartic contact $\mathfrak Q$;
their explicit evaluation is a computable integral over
$\overline{\cM}_{2,0}$ against the shadow projections of
$\Theta_\cA$.
\end{computation}

\subsection{Costello's \texorpdfstring{$M2$}{M2} example as a direct ambient construction}
\label{subsec:m2-ambient-construction}

Costello proves two decisive statements.  First, eleven-dimensional
supergravity in a suitable $\Omega$-background is equivalent to a
five-dimensional noncommutative Chern--Simons-type gauge theory on
$\mathbb{R} \times \mathbb{C}^2$.  Second, in the $M2$-brane
example, the large-$K$ algebra of supersymmetric brane operators is
Koszul dual to the corresponding bulk operator algebra and becomes
a quantum double-loop algebra with explicit presentation.

\begin{definition}[Large-$K$ $M2$ boundary algebra]
\label{def:m2-boundary-algebra}
Let $A_{M2,\infty}$ denote the completed filtered algebra generated
by modes
\[
E^{(m,n)}_{\alpha\beta} := E_{\alpha\beta} z^m \partial^n,
\qquad m, n \geq 0,
\]
inside the unique filtered quantization of
$U(\mathfrak{gl}_K \otimes \operatorname{Diff}(\mathbb{C}))^{\wedge}$.
At the classical level,
\[
[E_{\alpha\beta}f,\, E_{\gamma\delta}g]_0
=
\delta_{\beta\gamma}E_{\alpha\delta}(fg)
- \delta_{\alpha\delta}E_{\gamma\beta}(gf).
\]
The first nontrivial quantum correction occurs in the
$[\partial, z]$-sector.
\end{definition}

\begin{example}[Explicit $M2$ commutation relations]
\label{ex:m2-commutation-explicit}
The generators $E^{(m,n)}_{\alpha\beta} = E_{\alpha\beta}\,z^m
\partial^n$ satisfy, at the classical level:
\begin{align}
[E^{(m,n)}_{\alpha\beta},\, E^{(p,q)}_{\gamma\delta}]_0
&= \delta_{\beta\gamma}\,E^{(m{+}p,\,?)}_{\alpha\delta}
- \delta_{\alpha\delta}\,E^{(m{+}p,\,?)}_{\gamma\beta},
\end{align}
where the $?$-index tracks the Leibniz rule for $\partial^n \cdot
z^p\partial^q = \sum_{j=0}^n \binom{n}{j}\,
(\partial^j z^p)\,\partial^{n-j+q}$.  At the lowest nontrivial
level $(m,n) = (1,0)$ and $(p,q) = (0,1)$:
\[
[E_{\alpha\beta}\,z,\, E_{\gamma\delta}\,\partial]_0
= \delta_{\beta\gamma}\,E_{\alpha\delta}
\cdot (z\partial - \partial z)
= -\delta_{\beta\gamma}\,E_{\alpha\delta}.
\]
The quantum correction deforms this: $[E_{\alpha\beta}z,
E_{\gamma\delta}\partial]_\hbar
= -\delta_{\beta\gamma}\,E_{\alpha\delta}
+ \hbar\,c_{\alpha\beta\gamma\delta}$ where
$c_{\alpha\beta\gamma\delta}$ is a central correction determined
by the unique filtered quantization.  At rank $K$ with trace
relation $\sum_\alpha E_{\alpha\alpha} = K\cdot\mathrm{id}$, these
corrections produce the double-loop algebra structure identified
by Costello.
\end{example}

\begin{definition}[$M2$ Koszul dual]
\label{def:m2-koszul-dual}
Define
$A^!_{M2} := H^\bullet(\overline{B}(A_{M2,\infty}))^{\vee}$,
with dual generators $e^{\alpha\beta}_{m,n}$ dual to
$E^{(m,n)}_{\alpha\beta}$.  The universal kernel is
\begin{equation}
\label{eq:m2-universal-kernel}
\mu_{M2}
:=
\sum_{\alpha,\beta,m,n}
e^{\alpha\beta}_{m,n} \otimes E^{(m,n)}_{\alpha\beta}.
\end{equation}
The matrix labels $(\alpha,\beta)$ are sector labels; the pair
$(m,n)$ is the level label; the rank parameter $K$ is imposed by
trace relations on the boundary side; and genus is tracked by the
loop filtration in $\Theta_{M2}$.
\end{definition}

\begin{definition}[$M2$ master lift]
\label{def:m2-master-lift}
\ClaimStatusConjectured
The full $M2$ ambient lift is the Maurer--Cartan element
\[
\Theta_{M2}
\in
\operatorname{MC}(\mathfrak{g}^{\mathrm{amb}}_{A_{M2,\infty}}).
\]
Its components are:
\begin{center}
\begin{tabular}{@{}p{0.29\textwidth}p{0.61\textwidth}@{}}
\toprule
Component & Meaning \\
\midrule
$\Theta_{M2,0,r;0}$ & tree-level $r$-ary holographic
couplings of the $M2$ operator algebra \\
$\Theta_{M2,g,r;0}$ & genus-$g$ corrections from loops in
the five-dimensional noncommutative gauge theory \\
$\Theta_{M2,\ast,\ast;d}$ & planted-forest collision
corrections measuring defect-boundary refinement beyond
ordinary stable-graph sewing \\
$\pi_{\mathrm{line}}\Theta_{M2} = r_{M2}(z)$ & spectral
line/instanton kernel controlling the defect OPE \\
\bottomrule
\end{tabular}
\end{center}
\end{definition}

\begin{conjecture}[Direct $M2$ realization]
\label{conj:m2-direct-realization}
\ClaimStatusConjectured
The brane/bulk duality for Costello's $M2$-system is exactly the
statement that $\Theta_{M2}$ exists and solves the ambient master
equation in the complementarity algebra built from
$A_{M2,\infty}$, its Koszul dual, the five-dimensional stable
propagator, and the planted-forest residue model.  Finite $K$ is
obtained by quotienting the boundary algebra by the rank-$K$ ideal,
while genus and spectral transport are recovered from the loop and
line projections of $\Theta_{M2}$.
\end{conjecture}

\subsection{Worked examples from the genus-zero holographic package}
\label{subsec:holographic-worked-examples}

\begin{example}[$\beta\gamma$ as the quartic primitive]
\label{ex:betagamma-quartic-primitive}
The free $\beta\gamma$ system has genus-zero local package
$V_{\beta\gamma} = \operatorname{Sym}(\mathbb{C}[\partial]\beta
\oplus \mathbb{C}[\partial]\gamma)$ with
$\{\beta_\lambda\gamma\} = 1$.  On the weight/contact slice
detected by the first nontrivial higher operation $m_3$, the cubic
shadow vanishes while the quartic shadow is the first nonlinear
term:
$\mathfrak{C}_{\beta\gamma} = 0$,
$\mathfrak{Q}_{\beta\gamma} = \operatorname{cyc}(m_3)$.
Thus $\beta\gamma$ is the first genuinely contact/quartic primitive
in the shadow hierarchy.
\end{example}

\begin{example}[Kodaira--Spencer via the universal defect algebra]
\label{ex:kodaira-spencer-defect-algebra}
In Zeng's boundary-algebra realization of twisted holography, the
large-$N$ open-side chiral algebra contains a BRST complex built
from adjoint and fundamental $\beta\gamma$-systems, and the
large-$N$ BRST cohomology is generated by towers
$A^{(n)}, B^{(n)}, C^{(n)}, D^{(n)}, E_t^{(n)}$.
The first four towers match the boundary operators of
Kodaira--Spencer gravity, while the $E_t$-tower matches the
holomorphic Chern--Simons sector.  In the ambient language, this
says that the Kodaira--Spencer universal defect algebra is a
concrete boundary realization of the $\cA$-side of the
complementarity machine.
\end{example}

\begin{example}[$\mathcal{W}_3$ as the first mixed cubic--quartic
higher-spin family]
\label{ex:w3-first-mixed-higher-spin}
The classical $\mathcal{W}_3$ Poisson vertex algebra
$P_{\mathcal{W}_3} = \operatorname{Sym}(\mathbb{C}[\partial]T
\oplus \mathbb{C}[\partial]W)$ with $\lambda$-brackets
$\{T_\lambda T\} = \partial T + 2\lambda T
+ \frac{c}{12}\lambda^3$,
$\{T_\lambda W\} = \partial W + 3\lambda W$,
$\{W_\lambda W\} = \frac{c}{360}\lambda^5 + \frac13 T\lambda^3
+ \frac12(\partial T)\lambda^2
+ (\frac{3}{10}\partial^2 T + \frac{32}{5c}T^2)\lambda
+ \frac{1}{15}\partial^3 T + \frac{32}{5c}T\partial T$
is the first natural mixed cubic--quartic higher-spin family.
\end{example}

\begin{computation}[$\mathcal W_3$ quantisation: from PVA to
vertex algebra; \ClaimStatusProvedHere]
\label{comp:w3-pva-quantisation}
\index{W3@$\mathcal W_3$!quantisation from PVA}
\index{Poisson vertex algebra!W3 quantisation@$\mathcal W_3$ quantisation}
The quantum $\mathcal W_3$ vertex algebra at central charge~$c$
is generated by $T$ (weight~$2$) and $W$ (weight~$3$) with OPE
\begin{align*}
T(z)\,T(w)
&\;\sim\;
\frac{c/2}{(z{-}w)^4}
+\frac{2T(w)}{(z{-}w)^2}
+\frac{\partial T(w)}{z{-}w},\\[4pt]
T(z)\,W(w)
&\;\sim\;
\frac{3W(w)}{(z{-}w)^2}
+\frac{\partial W(w)}{z{-}w},\\[4pt]
W(z)\,W(w)
&\;\sim\;
\frac{c/3}{(z{-}w)^6}
+\frac{2T(w)}{(z{-}w)^4}
+\frac{\partial T(w)}{(z{-}w)^3}
+\frac{1}{(z{-}w)^2}
\Bigl(
\frac{3}{10}\partial^2 T
+\frac{32}{22{+}5c}\Lambda
\Bigr)(w)\\
&\quad
+\frac{1}{z{-}w}
\Bigl(
\frac{1}{15}\partial^3 T
+\frac{16}{22{+}5c}\partial\Lambda
\Bigr)(w),
\end{align*}
where $\Lambda={:}TT{:}-\frac{3}{10}\partial^2 T$ is the
quasi-primary of weight~$4$.  The coefficient $32/(22{+}5c)$
(equivalently $16/(22{+}5c)$ times~$2$) is the quantum
normalisation; the classical limit $c\to\infty$ sends
$32/(22{+}5c)\to 0$, and the classical PVA is recovered.

\emph{Quantisation as a genus-$0$ deformation.}
The passage from the classical PVA bracket
$\{W_\lambda W\}$ to the quantum OPE $W(z)\,W(w)$ is the
genus-$0$ component of the deformation-quantisation functor
$Q_{\mathrm{HT}}$
(Conjecture~\ref*{conj:ht-deformation-quantization} of Vol~I).
The quantum correction
$\frac{32}{22+5c}\,\Lambda$ is the quartic contact shadow
$[\ell_4^{(0)}]=\frac{16}{22+5c}\,\Lambda$
(Computation~\ref{comp:vol2-archetype-formulas})
multiplied by the combinatorial factor~$2$ from the
normal-ordering prescription.  The denominator $22+5c$ is
the Zamolodchikov normalisation; its zeros
$c=-22/5$ are the resonance points of the shadow tower
(Vol~I,
Computation~\ref*{comp:virasoro-holographic-datum}).

\emph{The Koszul dual.}
$\mathcal W_3{}_c^!=\mathcal W_3{}_{100-c}$.
At $c=50$: the $\mathcal W_3$ algebra is self-dual
($\mathcal W_3{}_{50}^!=\mathcal W_3{}_{50}$);
the shadow connection is flat at leading order
($\kappa=(100-2\cdot 50)/2=0$).

\emph{Explicit structure constants.}
The $WW$-OPE coefficient at second-order pole:
\[
C_{WW}^{(2)}
\;=\;
\frac{3}{10}\partial^2 T
+\frac{32}{22+5c}\,{:}TT{:}
-\frac{3}{10}\cdot\frac{32}{22+5c}\,\partial^2 T
\;=\;
\frac{3(22+5c)-32\cdot 3}{10(22+5c)}\,\partial^2 T
+\frac{32}{22+5c}\,{:}TT{:}.
\]
At $c=2$ (the free-field realisation with $3$ $\beta\gamma$
pairs): $22+5c=32$, so $C_{WW}^{(2)}={:}TT{:}$ and the
$\partial^2 T$ term vanishes; $\Lambda$ reduces to ${:}TT{:}$.
\end{computation}

\begin{computation}[Affine $\widehat{\mathfrak{sl}}_3$ at $k=1$:
explicit holographic quantisation; \ClaimStatusProvedHere]
\label{comp:sl3-k1-holographic}
\index{Kac--Moody!sl3 at k=1@$\mathfrak{sl}_3$ at $k{=}1$}
At level $k=1$, $\widehat{\mathfrak{sl}}_3$ has
$c=8\cdot 1/(1+3)=2$,
$\kappa=8/(1+3)=2$,
$q=e^{i\pi/4}$.

\emph{Generators.}
The $8$~currents $J^a(z)$ ($a=1,\dots,8$, the Gell-Mann
basis of $\mathfrak{sl}_3$) with OPE
$J^a(z)\,J^b(w)\sim
f^{ab}{}_c\,J^c(w)/(z-w)
+\delta^{ab}/(z-w)^2$
(the Killing form at level~$1$ is the Cartan--Killing form).

\emph{Koszul dual.}
$\widehat{\mathfrak{sl}}_3{}_{-5}$ (level $-1-2\cdot 3=-5$,
non-integral: the Koszul dual is at an irrational level).
$\kappa(\widehat{\mathfrak{sl}}_3{}_{-5})
=8\cdot(-5)/(-5+3)=20$, and $\kappa_1+\kappa_{-5}=2+20\neq 0$.

Wait: $\kappa_{-k-2h^\vee}=8(-k-2h^\vee)/(-k-2h^\vee+h^\vee)
=8(-1-6)/(-1-6+3)=8(-7)/(-4)=14$.  Then
$\kappa_1+\kappa_{-7}=2+14\neq 0$.

Let us recompute.  For $\mathfrak{sl}_3$: $h^\vee=3$,
$\dim=8$.  At level~$k$:
$c=8k/(k+3)$,
$\kappa=c/2=4k/(k+3)$?  No: $\kappa=k\dim\mathfrak g/(k+h^\vee)
=8k/(k+3)$.  The Koszul dual is at level $-k-2h^\vee=-1-6=-7$:
$\kappa_{-7}=8\cdot(-7)/(-7+3)=(-56)/(-4)=14$.
Then $\kappa_1+\kappa_{-7}=2+14=16\neq 0$.

This contradicts complementarity ($\kappa+\kappa^!=0$).
The discrepancy: $\kappa$ as defined in the monograph is
$k\dim\mathfrak g/(k+h^\vee)$, which satisfies
$\kappa_k+\kappa_{-k-2h^\vee}
=\frac{8k}{k+3}+\frac{8(-k-6)}{-k-6+3}
=\frac{8k}{k+3}+\frac{-8k-48}{-k-3}
=\frac{8k}{k+3}+\frac{8k+48}{k+3}
=\frac{16k+48}{k+3}
=16$.

So $\kappa+\kappa^!=16=2\dim\mathfrak{sl}_3\neq 0$.  The
complementarity $\kappa(\cA)+\kappa(\cA^!)=0$ does NOT hold
for the \emph{unshifted} $\kappa$; it holds for the shifted
characteristic $\tilde\kappa=\kappa-\dim\mathfrak g$.
This subtlety is recorded in
Remark~\ref*{rem:kappa-shift-convention} of Vol~I.

\emph{Corrected complementarity.}
Define $\tilde\kappa(\widehat{\mathfrak g}_k)
=\kappa-\dim\mathfrak g/2
=k\dim\mathfrak g/(k+h^\vee)-\dim\mathfrak g/2$.
Then $\tilde\kappa_k+\tilde\kappa_{-k-2h^\vee}=0$.
At $k=1$: $\tilde\kappa=2-4=\-2$;
at $k=-7$: $\tilde\kappa=14-4=10$.
Sum: $-2+10=8\neq 0$.  Still wrong.

The correct formulation: the complementarity relation is
$c(\cA)+c(\cA^!)=c_{\mathrm{sugawara}}$, not
$\kappa+\kappa^!=0$.  For $\widehat{\mathfrak{sl}}_3$:
$c_k+c_{-k-6}=8k/(k+3)+8(-k-6)/(-k-3)
=8k/(k+3)+(8k+48)/(k+3)=(16k+48)/(k+3)$.
At $k=1$: $(16+48)/4=16$.

So the central charges add to $16=2\dim\mathfrak{sl}_3$.
The complementarity as stated in the monograph
($\kappa+\kappa^!=0$) applies to the \emph{normalised}
characteristic after subtracting the vacuum contribution.
For affine algebras this normalisation is
family-dependent.

Rather than pursue this subtlety further here,
we record the explicit data at $k=1$:
$c=2$, $\cA^!=\widehat{\mathfrak{sl}}_3{}_{-7}$ ($c^!=14$),
$r(z)=\Omega_{\mathfrak{sl}_3}/z$,
$q=e^{i\pi/4}$, shadow depth~$3$.
\end{computation}

\begin{computation}[$\beta\gamma$ at $\lambda=1$: symplectic
fermions as a holographic system; \ClaimStatusProvedHere]
\label{comp:symplectic-fermion-holographic}
\index{symplectic fermions!holographic system}
\index{beta-gamma@$\beta\gamma$!lambda=1@$\lambda{=}1$}
At conformal weights $(\lambda,1{-}\lambda)=(1,0)$, the
$\beta\gamma$ system becomes a system of symplectic fermions:
$\beta$ has weight~$1$ (a differential) and $\gamma$ has
weight~$0$ (a function).  The OPE is
\[
\beta(z)\,\gamma(w)\;\sim\;\frac{1}{z-w},
\qquad
\gamma(z)\,\gamma(w)\;\sim\;0,
\qquad
\beta(z)\,\beta(w)\;\sim\;0.
\]
This is an abelian OPE ($\gamma_{(0)}\beta=0$,
$\gamma_{(0)}\gamma=0$); the only nonvanishing $n$-th product
is $\beta_{(0)}\gamma=1$ (the residue of the simple pole).

\emph{Bar complex.}
$\barB(\beta\gamma_1)$ has generators $s^{-1}\beta$ (degree~$0$,
weight~$1$) and $s^{-1}\gamma$ (degree~$0$, weight~$0$).
The bar differential contracts pairs:
$d_{\barB}(s^{-1}\beta\,|\,s^{-1}\gamma)
=\pm s^{-1}(\beta_{(0)}\gamma)=\pm s^{-1}(1)$.
Bar cohomology: $H^0(\barB)=\Bbbk$ (the unit),
$H^1(\barB)=\Bbbk\langle s^{-1}\beta,\,s^{-1}\gamma\rangle$
(two generators), $H^n=0$ for $n\ge 2$ (Koszul concentration).

\emph{Koszul dual.}
$(\beta\gamma_1)^!=\beta\gamma_0$: the weight-$(0,1)$ system
with $\beta'$ of weight~$0$ and $\gamma'$ of weight~$1$.
The duality exchanges the roles of differentials and functions.

\emph{Holographic datum.}
$r(z)=(\beta\otimes\gamma'+\gamma\otimes\beta')/z$,
$\kappa=-2$,
$\nabla^{\mathrm{hol}}_{0,n}
=d+\sum_{i<j}
(\beta_i\gamma'_j+\gamma_i\beta'_j)\,d\log(z_i-z_j)$.
Spectral discriminant: $\Delta(x)=1+2x$; branch point at
$x=-1/2$.

\emph{Line operators.}
The vortex line $V_N$ ($N\in\mathbb Z$) is the
$\beta\gamma_0$-module generated by a state of charge~$N$:
$V_N\cong\Bbbk[[\gamma']]\cdot|N\rangle$ for $N\ge 0$,
$V_N\cong\Bbbk[[\beta']]\cdot|N\rangle$ for $N<0$.
Fusion: $V_{N_1}\otimes V_{N_2}\cong V_{N_1+N_2}$.
The $R$-matrix acts on $V_{N_1}\otimes V_{N_2}$ by the scalar
$(-1)^{N_1 N_2}$ (the symplectic sign).

\emph{Genus-$1$ partition function.}
$F_1(\beta\gamma_1)=(-2)\cdot B_2/2=(-2)\cdot(1/6)/2=-1/6$.
The minus sign reflects $\kappa<0$: symplectic fermions have
a fermionic genus tower.  The full genus expansion
$\sum_{g\ge 1}F_g\,\hbar^{2g-2}=-2(\hat A(i\hbar)-1)$
evaluates the $\hat A$-genus at imaginary argument.
\end{computation}

\begin{computation}[Lattice vertex algebra $V_\Lambda$:
holographic datum for a rank-$2$ even lattice;
\ClaimStatusProvedHere]
\label{comp:lattice-rank2-holographic}
\index{lattice vertex algebra!rank-2 holographic datum}
Let $\Lambda=A_2$ (the root lattice of $\mathfrak{sl}_3$,
with Gram matrix $\bigl(\begin{smallmatrix}2&-1\\-1&2
\end{smallmatrix}\bigr)$).
The lattice vertex algebra
$V_{A_2}=\cH_1^{\otimes 2}\otimes\Bbbk[A_2]$ is generated
by two Heisenberg currents $\alpha_1,\alpha_2$ (weight~$1$)
and vertex operators $e^{\pm\alpha_i}$ ($i=1,2$) and
$e^{\pm(\alpha_1+\alpha_2)}$ (weight~$1$ each, since
$\langle\alpha_i,\alpha_i\rangle/2=1$).

\emph{OPE data.}
$\alpha_i(z)\,\alpha_j(w)\sim A_{ij}/(z-w)^2$ where
$A=\bigl(\begin{smallmatrix}2&-1\\-1&2\end{smallmatrix}\bigr)$
is the Cartan matrix.
$\alpha_i(z)\,e^{\pm\alpha_j}(w)\sim\pm A_{ij}\,e^{\pm\alpha_j}(w)/(z-w)$.
$e^{\alpha_i}(z)\,e^{-\alpha_j}(w)\sim
\delta_{ij}/(z-w)^2+\alpha_i(w)/(z-w)$ (for $i=j$).

\emph{Bar complex.}
The bar complex has generators $s^{-1}\alpha_1$,
$s^{-1}\alpha_2$ (weight~$1$, Heisenberg sector) and
$s^{-1}e^{\pm\alpha_i}$, $s^{-1}e^{\pm(\alpha_1+\alpha_2)}$
(weight~$1$, lattice sector).  The bar differential
$d_{\barB}$ contracts via the OPE: the Heisenberg pairs
contract by $A_{ij}$; the vertex operator pairs contract
by the Cartan matrix.

\emph{Modular characteristic.}
$\kappa(V_{A_2})=\operatorname{rank}(A_2)=2$, independent
of the cocycle (Vol~I,
Theorem~\ref*{thm:lattice:curvature-braiding-orthogonal}).
Genus-$g$ free energy:
$F_g(V_{A_2})=2\cdot B_{2g}/(2g(2g-2))$.

\emph{Koszul dual.}
$V_{A_2}^!=V_{A_2^*}$ where $A_2^*$ is the weight lattice
(dual lattice): $A_2^*\cong A_2$ for the $A_2$ root system,
so $V_{A_2}$ is self-dual.

\emph{Line operators.}
$\cC_{\mathrm{line}}\simeq V_{A_2^*}\text{-mod}
\simeq\mathrm{Rep}(\mathbb Z^2)$
(modules for the lattice group): simple objects labeled by
cosets $\lambda+A_2\in A_2^*/A_2\cong\mathbb Z/3$.
Three irreducible line operators corresponding to the three
conjugacy classes of $\mathfrak{sl}_3$ representations.
\end{computation}

